\listfiles
\documentclass{article}
\hbadness=10000 \newcount\hbadness \hfuzz=\maxdimen % Make TeX shut up.

%\usepackage{pkgindoc}
\makeatletter
\let\@onlypreamble\@gobble
\makeatother

\usepackage[T1]{fontenc}
\usepackage[utf8]{inputenc}

\usepackage{booktabs}
\usepackage{array,chngpage}
\usepackage[demo]{graphicx}

\usepackage{caption,float,longtable,subfig}
\usepackage{algorithm2e}

\newcommand*\bug[1]{%
  \clearpage\section*{Bug #1}\noindent
  \bugi}
\newcommand*\bugi[2][]{%
  \def\temp{#1}\ifx\temp\empty
    \textit{(solved in #2)}
  \else
    \textit{(introduced in #1, solved in #2)}
  \fi
  \par\noindent\ignorespaces}
\newenvironment{example}{\medskip\hrule\nopagebreak\captionsetup}{\nopagebreak\hrule}

\begin{document}

%******************************************************************************%

\bug{04-01-17}{v3.0a}

\begin{verbatim}
Date: Sat, 17 Jan 2004 12:24:25 +0100
From: Frank Mittelbach <frank.mittelbach@latex-project.org>
Reply-To: Frank Mittelbach <frank.mittelbach@latex-project.org>
To: Steven Douglas Cochran <Steven_Douglas_Cochran@maps.cs.cmu.edu>
Cc: Axel Sommerfeldt <axel@sommerfeldt.net>,
	David Carlisle <davidc@nag.co.uk>,
	Harald Harders <h.harders@tu-bs.de>
Subject: Re: Update: The subfig package 1.1

[...]
that would work except that it shows a bug in the caption package (sorry
Axel). The caption package doesn't allow

  \captionsetup[table]   {position=top,aboveskip=5pt}

ie spaces between the option argument and the main argument.
The reason is the following definition:

%FMi need to pass arg #2 or spaces might show up in def!
\def\caption@setuptype[#1]#2{%
  \@ifundefined{caption@typ@#1}%
%FMi missing percent
    {\@namedef{caption@typ@#1}{#2}}%
    {\expandafter\l@addto@macro\csname caption@typ@#1\endcsname,%
     \expandafter\l@addto@macro\csname caption@typ@#1\endcsname{#2}}}

the missing percent is not important (it only save space) but one has to pass
#2 otherwise \@namedef will actually define a macro

 \caption@typ@table 

that has to be followed by a space

Axel, can you please fix that since one should always be allowed to put spaces
in such a place without doing any harm

thanks
frank
\end{verbatim}

\begingroup
\captionsetup[table]   {position=top,aboveskip=5pt}
\DeclareCaptionLabelSeparator{period-newline}  {.\newline}
\endgroup

%******************************************************************************%

\bug{04-05-05}{v3.0b}
\bugi[v3.0j]{v3.0l}

\begin{verbatim}
Newsgroups: de.comp.text.tex
Subject: subfig und caption mit format=hang
Date: Wed, 5 May 2004 20:52:12 +0200
Message-ID: <s4db7c.pm3.ln@pjl.fqdn.th-h.de>

Hallo,

mal wieder ein kleines Problem:
Beim u.a. Minimalbeispiel werden nicht nur die Folgezeilen der Legende,
sondern bereits die erste Zeile _inklusive_ "Numerierung" eingezogen.
Sieht etwa so aus:

+----------+
|          |
|          |
|          |
|          |
|          |
|          |
|          |
|          |
|          |
|          |
+----------+

  (a) laaaaaaaaaaang
  kurz kurz
  kurz 

Mache ich "laaaaaaaang" kürzer, fängt die Legende schön linksbündig an.

Das Beispiel:

\documentclass{article}
\usepackage[format=hang]{caption}
\usepackage{subfig}
\listfiles
\begin{document}
\begin{figure}
\subfloat[laaaaaaaaaaaaaaaaang kurz kurz kurz]{\rule{2cm}{5cm}}
\end{figure}
\end{document}

Die fraglichen Pakete sind m.E. auf dem neusten Stand:
 article.cls    2001/04/21 v1.4e Standard LaTeX document class
  size10.clo    2001/04/21 v1.4e Standard LaTeX file (size option)
 caption.sty    2004/01/23 v3.0a Customising captions (AS)
  keyval.sty    1999/03/16 v1.13 key=value parser (DPC)
  subfig.sty    2004/01/28 ver: 1.2 subfig package
ragged2e.sty    1999/06/08 v1.02 Ragged2e Package (MS)
everysel.sty    1999/06/08 v1.03 EverySelectfont Package (MS)

Weiß jemand Abhilfe?

Danke,
Peter
\end{verbatim}

\begin{example}{format=hang}
%\captionsetup[figure]{twoside=false}

\begin{figure}[H]
\subfloat[laaaaaaaaaaaaaaaaang kurz kurz kurz]{\rule{2cm}{5cm}}
\end{figure}

\end{example}

%******************************************************************************%

\bug{04-05-06}{v3.0b}

\begin{verbatim}
Newsgroups: de.comp.text.tex
Subject: Re: subfig und caption mit format=hang
Date: Thu, 6 May 2004 14:24:26 +0200
Message-ID: <qpad7c.942.ln@pjl.fqdn.th-h.de>

Hallo Axel,

Axel Sommerfeldt <sommerfee@despammed.com> wrote:
> 
> Dann sollte es gehen. Magst du das mal ausprobieren und bestätigen? Das 
> wäre nett.

Das leistet das Verlangte. Allerdings gibt es noch folgende Unschönheit:
Im ersten Wort der Legende findet keine Silbentrennung statt, selbst
dann nicht, wenn man einen trennbaren Text wählt (das erste Bsp. war in
dieser Hinsicht suboptimal). Lasse ich "format=hang" weg, wird getrennt!

Das neue Beispiel:

\documentclass{article}
\usepackage{caption}
\usepackage[format=hang]{subfig}
\begin{document}
\begin{figure}
\subfloat[Donaudampfschifffahrtsgesellschaft]{\rule{2cm}{5cm}}
\end{figure}
\end{document}

Peter
\end{verbatim}

\begin{example}{format=hang}

\hyphenation{Do-nau-dampf-schiff-fahrts-ge-sell-schaft}
\begin{figure}[H]
\subfloat[Donaudampfschifffahrtsgesellschaft]{\rule{2cm}{5cm}}
\end{figure}
%\begin{figure}[H]
%\subfloat[Donaudampfschifffahrtsgesellschaft Donaudampfschifffahrtsgesellschaft Donaudampfschifffahrtsgesellschaft]{\rule{2cm}{5cm}}
%\end{figure}

\end{example}

%******************************************************************************%

\bug{04-07-15}[v3.0b]{v3.0c}

\begin{verbatim}
Newsgroups: comp.text.tex
Subject: Compile error with TLC2 6-5-15.ltx
Date: Thu, 15 Jul 2004 09:42:47 -0700
Message-ID: <qlcdf01sk66hjm3da7kdi5mvip2htsv7r1@4ax.com>

The following statements are belong to the log file.
My caption.sty is made at 5/16/2004. Do you know what is wrong with
it?

! Undefined control sequence.
<argument> \SC@CAPsetup 
                        
l.33 \end{SCtable}
\end{verbatim}

%******************************************************************************%

\bug{04-08-04}{v3.0d}

\begin{verbatim}
Newsgroups: de.comp.text.tex
Date: Wed, 4 Aug 2004 08:15:02 +0200
Message-ID: <41107ee6$0$412$ba620e4c@news.skynet.be>
Subject: caption +longtable

1. Caption klebt an der Longtable (als Vergleich hab ich noch eine 
normale Tabelle mitgepostet)
2. Stilistische Frage: würdet ihr den Seitenumbruch in der Longtable 
anders gestalten?

Danke für die Mühe.

Package: scrlfile 2004/01/07 v2.9p LaTeX2e KOMA package
Package: booktabs 2003/03/28 v1.618 publication quality tables
Package: longtable 2004/02/01 v4.11 Multi-page Table package (DPC)
Package: caption 2004/05/16 v3.0b Customising captions (AS)


\documentclass[pdftex, headsepline]{scrbook}

\usepackage[ngerman]{babel}
\usepackage[applemac]{inputenc}
\usepackage[T1]{fontenc}	
\usepackage{booktabs}
\usepackage{longtable}
\usepackage[indention=1em,font=small,labelfont=bf,margin=1.5em,parindent=1em]{caption}
\newcommand{\ebf}[1]{{\emph{\textbf{#1}}}}

\begin{document}

\begin{table}[htbp]
\setlength{\tabcolsep}{0.3em}
\footnotesize
\centering
\begin{tabular}{@{} *{11}{l} @{}}
	\toprule
	\multicolumn{11}{@{} l @{}}{\ebf{Ziel:} Happy TeXing} \\
	\midrule
	\ebf{OR}	& 1.	& \multicolumn{9}{l@{}}{de.comp.text.tex} \\
			&	& \ebf{OR}	& 1.	& \multicolumn{7}{l@{}}{de.comp.text.tex} \\
			&	& 			&	& \ebf{OR}	& 1.	& \multicolumn{5}{l@{}}{de.comp.text.tex} \\
			&	& 			&	&			& 	& \ebf{OR}	& 1.	& \multicolumn{3}{l@{}}{de.comp.text.tex} \\
			&	& 			&	&			& 	&			& 2.	& \multicolumn{3}{l@{}}{de.comp.text.tex} \\
			&	&			& 	& 			& 2.	& \multicolumn{5}{l@{}}{de.comp.text.tex} \\
			&	& 			&	&			& 	& \ebf{OR}	& 1.	& \multicolumn{3}{l@{}}{de.comp.text.tex} \\
			&	& 			&	&			& 	&			& 2.	& \multicolumn{3}{l@{}}{de.comp.text.tex} \\
			&	& 			&	&			& 	& 			&	& \ebf{AND}	& 1.	& 
\multicolumn{1}{l@{}}{de.comp.text.tex} \\
			&	& 			&	&			& 	&			& 	&			& 2.	& \multicolumn{1}{l@{}}{de.comp.text.tex} \\
	\bottomrule
\end{tabular}
\caption{Happy TeXing: de.comp.text.tex}
\label{tab:Angriffsbaum-A}
\end{table}
Hier ist ein normaler Abstand zur Tabelle.


{%
\setlength{\tabcolsep}{0.3em}
\footnotesize
\centering
\begin{longtable}{@{} *{11}{l} @{}}
	\toprule
	\multicolumn{11}{@{} l @{}}{\ebf{Ziel:} Happy TeXing} \\
	\midrule
	\endfirsthead
	
	\midrule
	\endhead
	
	\midrule
	\multicolumn{11}{@{} r @{}}{Fortsetzung auf der nächsten Seite} \\
	\midrule
	\endfoot
	
	\bottomrule
	\caption{Happy TeXing: de.comp.text.tex}
	\label{tab:Angriffsbaum-B}
	\endlastfoot	
	\ebf{OR}	& 1.	& \multicolumn{9}{l@{}}{de.comp.text.tex} \\
			&	& \ebf{OR}	& 1.	& \multicolumn{7}{l@{}}{de.comp.text.tex} \\
			&	& 			&	& \ebf{OR}	& 1.	& \multicolumn{5}{l@{}}{de.comp.text.tex} \\
			&	& 			&	&			& 	& \ebf{OR}	& 1.	& \multicolumn{3}{l@{}}{de.comp.text.tex} \\
			&	& 			&	&			& 	&			& 2.	& \multicolumn{3}{l@{}}{de.comp.text.tex} \\
			&	&			& 	& 			& 2.	& \multicolumn{5}{l@{}}{de.comp.text.tex} \\
			&	& 			&	&			& 	& \ebf{OR}	& 1.	& \multicolumn{3}{l@{}}{de.comp.text.tex} \\
			&	& 			&	&			& 	&			& 2.	& \multicolumn{3}{l@{}}{de.comp.text.tex} \\
			&	& 			&	&			& 	& 			&	& \ebf{AND}	& 1.	& 
\multicolumn{1}{l@{}}{de.comp.text.tex} \\
			&	& 			&	&			& 	&			& 	&			& 2.	& \multicolumn{1}{l@{}}{de.comp.text.tex} \\
\end{longtable}
}
Hier ist kein normaler Abstand zur Tabelle und die Caption klebt an der 
Tabelle.

\end{document}
\end{verbatim}

\begin{example}{indention=1em,font=small,labelfont=bf,margin=1.5em,parindent=1em}
\captionsetup[table]{position=bottom}
\newcommand{\ebf}[1]{{\emph{\textbf{#1}}}}

\begin{table}[H]
\setlength{\tabcolsep}{0.3em}
\footnotesize
\centering
\begin{tabular}{@{} *{11}{l} @{}}
	\toprule
	\multicolumn{11}{@{} l @{}}{\ebf{Ziel:} Happy TeXing} \\
	\midrule
	\ebf{OR}	& 1.	& \multicolumn{9}{l@{}}{de.comp.text.tex} \\
			&	& \ebf{OR}	& 1.	& \multicolumn{7}{l@{}}{de.comp.text.tex} \\
			&	& 			&	& \ebf{OR}	& 1.	& \multicolumn{5}{l@{}}{de.comp.text.tex} \\
			&	& 			&	&			& 	& \ebf{OR}	& 1.	& \multicolumn{3}{l@{}}{de.comp.text.tex} \\
			&	& 			&	&			& 	&			& 2.	& \multicolumn{3}{l@{}}{de.comp.text.tex} \\
			&	&			& 	& 			& 2.	& \multicolumn{5}{l@{}}{de.comp.text.tex} \\
			&	& 			&	&			& 	& \ebf{OR}	& 1.	& \multicolumn{3}{l@{}}{de.comp.text.tex} \\
			&	& 			&	&			& 	&			& 2.	& \multicolumn{3}{l@{}}{de.comp.text.tex} \\
			&	& 			&	&			& 	& 			&	& \ebf{AND}	& 1.	& 
\multicolumn{1}{l@{}}{de.comp.text.tex} \\
			&	& 			&	&			& 	&			& 	&			& 2.	& \multicolumn{1}{l@{}}{de.comp.text.tex} \\
	\bottomrule
\end{tabular}
\caption{Happy TeXing: de.comp.text.tex}
\label{tab:Angriffsbaum-A}
\end{table}
Hier ist ein normaler Abstand zur Tabelle.


{%
\setlength{\tabcolsep}{0.3em}
\footnotesize
\centering
\begin{longtable}{@{} *{11}{l} @{}}
	\toprule
	\multicolumn{11}{@{} l @{}}{\ebf{Ziel:} Happy TeXing} \\
	\midrule
	\endfirsthead
	
	\midrule
	\endhead
	
	\midrule
	\multicolumn{11}{@{} r @{}}{Fortsetzung auf der nächsten Seite} \\
	\midrule
	\endfoot
	
	\bottomrule
	\caption{Happy TeXing: de.comp.text.tex}
	\label{tab:Angriffsbaum-B}
	\endlastfoot	
	\ebf{OR}	& 1.	& \multicolumn{9}{l@{}}{de.comp.text.tex} \\
			&	& \ebf{OR}	& 1.	& \multicolumn{7}{l@{}}{de.comp.text.tex} \\
			&	& 			&	& \ebf{OR}	& 1.	& \multicolumn{5}{l@{}}{de.comp.text.tex} \\
			&	& 			&	&			& 	& \ebf{OR}	& 1.	& \multicolumn{3}{l@{}}{de.comp.text.tex} \\
			&	& 			&	&			& 	&			& 2.	& \multicolumn{3}{l@{}}{de.comp.text.tex} \\
			&	&			& 	& 			& 2.	& \multicolumn{5}{l@{}}{de.comp.text.tex} \\
			&	& 			&	&			& 	& \ebf{OR}	& 1.	& \multicolumn{3}{l@{}}{de.comp.text.tex} \\
			&	& 			&	&			& 	&			& 2.	& \multicolumn{3}{l@{}}{de.comp.text.tex} \\
			&	& 			&	&			& 	& 			&	& \ebf{AND}	& 1.	& 
%\multicolumn{1}{l@{}}{de.comp.text.tex} \\
%			&	& 			&	&			& 	&			& 	&			& 2.	& \multicolumn{1}{l@{}}{de.comp.text.tex} \\
%			& 2.	& \multicolumn{9}{l@{}}{de.comp.text.tex} \\
%			&	& \ebf{OR}	& 1.	& \multicolumn{7}{l@{}}{de.comp.text.tex} \\
%			&	& 			&	& \ebf{OR}	& 1.	& \multicolumn{5}{l@{}}{de.comp.text.tex} \\
%			&	& 			&	&			& 	& \ebf{OR}	& 1.	& \multicolumn{3}{l@{}}{de.comp.text.tex} \\
%			&	& 			&	&			& 	&			& 2.	& \multicolumn{3}{l@{}}{de.comp.text.tex} \\
%			&	&			& 	& 			& 2.	& \multicolumn{5}{l@{}}{de.comp.text.tex} \\
%			&	& 			&	&			& 	& \ebf{OR}	& 1.	& \multicolumn{3}{l@{}}{de.comp.text.tex} \\
%			&	& 			&	&			& 	&			& 2.	& \multicolumn{3}{l@{}}{de.comp.text.tex} \\
%			&	& 			&	&			& 	& 			&	& \ebf{AND}	& 1.	& 
%\multicolumn{1}{l@{}}{de.comp.text.tex} \\
%			&	& 			&	&			& 	&			& 	&			& 2.	& \multicolumn{1}{l@{}}{de.comp.text.tex} \\
%			& 3.	& \multicolumn{9}{l@{}}{de.comp.text.tex} \\
%			&	& \ebf{OR}	& 1.	& \multicolumn{7}{l@{}}{de.comp.text.tex} \\
%			&	& 			&	& \ebf{OR}	& 1.	& \multicolumn{5}{l@{}}{de.comp.text.tex} \\
%			&	& 			&	&			& 	& \ebf{OR}	& 1.	& \multicolumn{3}{l@{}}{de.comp.text.tex} \\
%			&	& 			&	&			& 	&			& 2.	& \multicolumn{3}{l@{}}{de.comp.text.tex} \\
%			&	&			& 	& 			& 2.	& \multicolumn{5}{l@{}}{de.comp.text.tex} \\
%			&	& 			&	&			& 	& \ebf{OR}	& 1.	& \multicolumn{3}{l@{}}{de.comp.text.tex} \\
%			&	& 			&	&			& 	&			& 2.	& \multicolumn{3}{l@{}}{de.comp.text.tex} \\
%			&	& 			&	&			& 	& 			&	& \ebf{AND}	& 1.	& 
%\multicolumn{1}{l@{}}{de.comp.text.tex} \\
%			&	& 			&	&			& 	&			& 	&			& 2.	& \multicolumn{1}{l@{}}{de.comp.text.tex} \\
%			& 4.	& \multicolumn{9}{l@{}}{de.comp.text.tex} \\
%			&	& \ebf{OR}	& 1.	& \multicolumn{7}{l@{}}{de.comp.text.tex} \\
%			&	& 			&	& \ebf{OR}	& 1.	& \multicolumn{5}{l@{}}{de.comp.text.tex} \\
%			&	& 			&	&			& 	& \ebf{OR}	& 1.	& \multicolumn{3}{l@{}}{de.comp.text.tex} \\
%			&	& 			&	&			& 	&			& 2.	& \multicolumn{3}{l@{}}{de.comp.text.tex} \\
%			&	&			& 	& 			& 2.	& \multicolumn{5}{l@{}}{de.comp.text.tex} \\
%			&	& 			&	&			& 	& \ebf{OR}	& 1.	& \multicolumn{3}{l@{}}{de.comp.text.tex} \\
%			&	& 			&	&			& 	&			& 2.	& \multicolumn{3}{l@{}}{de.comp.text.tex} \\
%			&	& 			&	&			& 	& 			&	& \ebf{AND}	& 1.	& 
\multicolumn{1}{l@{}}{de.comp.text.tex} \\
			&	& 			&	&			& 	&			& 	&			& 2.	& \multicolumn{1}{l@{}}{de.comp.text.tex} \\
\end{longtable}
}
Hier ist kein normaler Abstand zur Tabelle und die Caption klebt an der 
Tabelle.

\end{example}

%******************************************************************************%

\bug{04-10-26}{v3.0d \& v3.0i}

\begin{verbatim}
Newsgroups: de.comp.text.tex
Subject: Re: subfig inkompatibel mit Komascript ?
Date: Tue, 26 Oct 2004 09:37:12 +0200
Message-ID: <clkupd$23m$04$1@news.t-online.com>

Axel Sommerfeldt wrote:

> Matthias Pospiech <matthiasPUNKTpospiech@gmx.de> wrote:
> 
>> Leider setzt caption bei Raggedright die erste Zeile nicht mehr mittig
>> (so wie KOMAscript) sondern auch nach Rechts. Ich sehe auch nicht wie
>> ich das caption abgewöhnen könnte.
> 
> Mir ist nicht ganz klar, was du damit meinst. Hast du ein KOMA-
> Minimalbeispiel für mich, was das produziert, wie du es haben möchtest?

Hier ist eins. Es ist nicht wirklich minimal, da ich auch selbst definierte
Befehle brauche, aber es zeigt sehr deutich den Unterschied. Für Komascript
Anweisungen einfach bei diesen die Kommentare wegnehmen.
----------------------------------------

\documentclass[a4paper, 11pt,german]{scrreprt}

\usepackage[NewCommands]{ragged2e}

%\setkomafont{caption}{\footnotesize\sffamily\RaggedRight}
% Schrift für
%\setkomafont{captionlabel}{\small}   % Schrift für 'Abbildung' usw.
%\setcapindent{0em}

\usepackage[margin=0pt,textfont={sf,footnotesize},labelfont={sf,small},%
  justification=RaggedRight,singlelinecheck=true]{caption}

\usepackage[final]{graphicx}
\usepackage{array}

\makeatletter
\newlength{\figcaptionwidth}
%
\newcommand\figcaption{\def\@captype{figure}\caption}
\newcommand\tabcaption{\def\@captype{table}\caption}
%
\newcommand\figurecaption[2]{
\setlength{\figcaptionwidth}{#1}
\\\begin{minipage}{\figcaptionwidth}
\figcaption{#2}
\end{minipage}
}
\makeatother

\newcolumntype{v}[1]{%
        >{\raggedright\hspace{0pt}}p{#1}%
}
\newenvironment{wide}[2]{%
\begin{list}{}{%
\setlength{\topsep}{0pt}%
\setlength{\leftmargin}{#1}%
\setlength{\rightmargin}{#2}%
\setlength{\listparindent}{\parindent}%
\setlength{\itemindent}{\parskip}%
}%
\item[]}%
{\end{list}}%

\newlength{\marginwidth}
\setlength{\marginwidth}{\marginparwidth}
\addtolength{\marginwidth}{\marginparsep}
        
\begin{document}

\begin{wide}{-\marginwidth}{0pt}
        %\tablestylecommon
        %\tablealtcolored
        \begin{flushright}
                        \begin{tabular}{*{7}{v{0.14\textwidth}}}
                \hline
 Current &
 Power LD A &
 Power LD B &
 LD-Power together &
 LD-Power behind pump optics &
 Losses through pump optics &
 IR Power Mephisto Module \tabularnewline
                \hline
\end{tabular}
        \end{flushright}
\tabcaption{\label{tab:leistung}Leistungsmessdaten eines Produktionslasers}
\end{wide}

\begin{center}
%\includegraphics[scale=1]{data/powerP100}
\fbox{Bild}
\figcaption{\label{fig:leistung}Pump und Laserausgangsleistung}
\end{center}

\begin{center}
        %\includegraphics[width=0.7\textwidth]{images/pdf/FPmeasure}
        \fbox{Bild}
        \figurecaption{0.7\textwidth}{\label{fig:FPmeasure} Emissionsspektrum eines
Mephistolasers aufgenommen mit einem Scanning Fabry-Perot Interferometer
(FPI)}
\end{center}


\end{document}
\end{verbatim}

\begin{example}{margin=0pt,textfont={sf,footnotesize},labelfont={sf,small},%
  justification=RaggedRight,singlelinecheck=true}

\newenvironment{wide}[2]{%
\begin{list}{}{%
\setlength{\topsep}{0pt}%
\setlength{\leftmargin}{#1}%
\setlength{\rightmargin}{#2}%
\setlength{\listparindent}{\parindent}%
\setlength{\itemindent}{\parskip}%
}%
\item[]}%
{\end{list}}%

\newcolumntype{v}[1]{%
        >{\raggedright\hspace{0pt}}p{#1}%
}

\newlength{\marginwidth}
\setlength{\marginwidth}{\marginparwidth}
\addtolength{\marginwidth}{\marginparsep}

\begin{wide}{-\marginwidth}{0pt}
        %\tablestylecommon
        %\tablealtcolored
        \begin{flushright}
                        \begin{tabular}{*{7}{v{0.14\textwidth}}}
                \hline
 Current &
 Power LD A &
 Power LD B &
 LD-Power together &
 LD-Power behind pump optics &
 Losses through pump optics &
 IR Power Mephisto Module \tabularnewline
                \hline
\end{tabular}
        \end{flushright}
\captionof{table}{\label{tab:leistung}Leistungsmessdaten eines Produktionslasers}
\end{wide}

\begin{wide}{-\marginwidth}{0pt}
        %\tablestylecommon
        %\tablealtcolored
        \begin{flushright}
                        \begin{tabular}{*{7}{v{0.14\textwidth}}}
                \hline
 Current &
 Power LD A &
 Power LD B &
 LD-Power together &
 LD-Power behind pump optics &
 Losses through pump optics &
 IR Power Mephisto Module \tabularnewline
                \hline
\end{tabular}
        \end{flushright}
\captionof{table}{Leistungsmessdaten eines Produktionslasers
 -- Leistungsmessdaten eines Produktionslasers -- Leistungsmessdaten eines Produktionslasers
 -- Leistungsmessdaten eines Produktionslasers}
\end{wide}

\end{example}

%******************************************************************************%

\bug{04-12-19}{v3.0d}

\begin{verbatim}
Newsgroups: de.comp.text.tex
Subject: Zeilenabstände mit \usepackage[font=small]{caption}
Date: Sun, 19 Dec 2004 14:59:42 +0100
Message-ID: <cq41in$qdn$03$1@news.t-online.com>

Hallo NG,

folgendes Minimalbeispiel

\documentclass[12pt,a4paper]{scrartcl}
\usepackage{ngerman}
\usepackage[latin1]{inputenc}
\usepackage[font=small]{caption}[2004/07/16]

\begin{document}
\begin{figure}
  \centering
  \caption{Dies ist ein Text, der nur zu Testzwecken eingefügt ist, um die
    Zeilenabstände sichtbar zu machen.  Hier taucht jetzt ein \AA{}
    auf, an dem aufgezeigt werden soll, dass eben die Zeilenabstände nicht
    stimme.  Und noch mehr Text taucht hier auf, damit auch gesichert
    mindestens vier Zeilen in dem fertigen Dokument stehen.}
\end{figure}
\end{document}

führt bei mir zu einem Dokument, in dem die Zeilenabstände nicht
gleichverteilt erscheinen.  Irgendeine Idee, woran das liegen
könnte?


Danke im voraus,
Jürgen


Newsgroups: de.comp.text.tex
Subject: Re: Zeilenabstände mit \usepackage[font=small]{caption}
Date: Sun, 19 Dec 2004 18:57:22 +0100
Message-ID: <3sebs0hkmsjobnn6di971hbqqpaiat9nd1@4ax.com>

On Sun, 19 Dec 2004 14:59:42 +0100, Juergen Wieferink wrote:

>Hallo NG,
>
>folgendes Minimalbeispiel

[snip]

>führt bei mir zu einem Dokument, in dem die Zeilenabstände nicht
>gleichverteilt erscheinen.  Irgendeine Idee, woran das liegen
>könnte?

Vermutlich an der Stütze die vor dem Label eingefügt wird. Wie man das
Ganze sauber löst kann ich dir nicht sagen, aber vielleicht reicht dir
unten angeführte Lösung. Jedenfalls kannst du den Unterschied sehen.

\usepackage{caption}

\DeclareCaptionFont{small*}{\small\setbox\strutbox=\box255}
\captionsetup{font=small*}

Wenn es sich nur um einzelne Zeichen handelt, kannst du diese auch in
eine Box geben, anschließend die Höhe der Box reduzieren und diese
ausgeben.

Prinzip:

\setlength{\fboxsep}{0pt}
\fbox{xxx\setbox0\hbox{\AA}\ht0=0pt\box0xxx}


Chris
\end{verbatim}

\begin{example}{font=small}

\begin{figure}[H]
  \centering
  \caption{Dies ist ein Text, der nur zu Testzwecken eingefügt ist, um die
    Zeilenabstände sichtbar zu machen.  Hier taucht jetzt ein \AA{}
    auf, an dem aufgezeigt werden soll, dass eben die Zeilenabstände nicht
    stimme.  Und noch mehr Text taucht hier auf, damit auch gesichert
    mindestens vier Zeilen in dem fertigen Dokument stehen.}
\end{figure}

\end{example}

%******************************************************************************%

\bug{05-01-23}{v3.0d}

\begin{verbatim}
Message-ID: <000701c5015d$04dc9b40$b9653ac2@home196ib52sok>
To: "Axel Sommerfeldt" <axel.sommerfeldt@arcor.de>
Subject: Re: rotfloat question
Date: Sun, 23 Jan 2005 18:04:27 +0300

[...]

If I build row of floats where first caption has last line with letters,
like g, j, p, q, y (with descenders), but in second caption there are not
any of them, so the baselines of that two captions are not aligned. I found
that you start text in caption@make stuff with strut. I think it is good to
finish text of caption with \strut too, or, like in footnotes, with
\@finalstrut\strutbox (this two-command combination adds only depth of
strut). [I redefined \caption@@@make in test file and get correct result.]

This \strut could get empty line only if someone writes like:
\caption{Text of caption

}
I never see such variant, and, in anyway, this is wrong.(?)

That not visible with usage of standard
\abovecaptionskip(\belowcaptionskip)=10pt, but if they reduced to
somewhat=3pt
the difference could be visible. This is also could be visible in beside
subcaptions.

Regards,
Olga
\end{verbatim}

%******************************************************************************%

\bug{05-04-08}{v3.0e}

\begin{verbatim}
Newsgroups: de.comp.text.tex
Subject: algorithm2e und caption
Date: Fri, 08 Apr 2005 10:35:51 +0200
Message-ID: <7cshi2-ped.ln1@ID-80222.user.uni-berlin.de>

Hallo,

ich wollte algorithm2e zusammen mit dem caption Paket benutzen. Leider
werden nach dem Laden von caption die Label der Algorithmen nicht mehr
gezeigt. Anstatt "Algorithm 5: Move" sehe ich nur noch "5: Move".

Kennt jemand das Problem


Newsgroups: de.comp.text.tex
Subject: Re: algorithm2e und caption
Date: Fri, 08 Apr 2005 11:08:22 +0200
Message-ID: <79uhi2-5ee.ln1@ID-80222.user.uni-berlin.de>

Alois Steindl wrote:

> Christoph Bartoschek writes:
> 
>> Hallo,
>> 
>> ich wollte algorithm2e zusammen mit dem caption Paket benutzen. Leider
>> werden nach dem Laden von caption die Label der Algorithmen nicht mehr
>> gezeigt. Anstatt "Algorithm 5: Move" sehe ich nur noch "5: Move".
>> 
>> Kennt jemand das Problem
>> 
> Wie wäre es mit einem Minimalbeispiel? Caption lässt sich gut
> konfigurieren. Eventuell Reihenfolge beim Laden der Pakete ändern.
> Alois

Hier ein relativ kleines Beispiel:


\documentclass{article}

\usepackage{algorithm2e}
%\usepackage{caption}

\begin{document}

\begin{algorithm}
$c \leftarrow d$ \;
\caption{Testalgorithm}
\end{algorithm}

\end{document}

Wenn man den Kommentar bei caption entfernt, verschwindet das Wort Algorithm
und die 1 ist auch nicht mehr fett.
Ein Vertauschen der Pakete bringt nichts. 
\end{verbatim}

\begin{example}{}
\begin{algorithm}
$c \leftarrow d$ \;
\caption{Testalgorithm}
\end{algorithm}
\end{example}

%******************************************************************************%

\bug{05-04-28}[v3.0d]{v3.0e}

\begin{verbatim}
Message-ID: <4270EAF4.700@web.de>
Date: Thu, 28 Apr 2005 15:53:56 +0200
To: caption@sommerfeldt.net
Subject: [caption] Fehler bei indent

Hallo Herr Sommerfeldt,

ich habe gerade ein Update des caption Paketes eingespielt und habe nun
folgendes Problem:

Die Option "indention" erzeugt nicht mehr einen Einzug für alle Zeilen
der Beschriftung, sondern einen "Auszug" der ersten Zeile. Das bedeutet,
dass die erste Zeile in den Seitenrand hineinragt. Ich verwende die
Paketeversionen

caption.sty    2004/11/28 v3.0d Customising captions (AS)
caption3.sty    2005/02/12 v3.0d caption3 kernel (AS)

und z.B. den Befehl

\usepackage[labelfont=bf,indention=2cm]{caption}

Zur Demonstration folgt eine Beisppieldatei:

--------------------
\documentclass[12pt]{article}

\usepackage[labelfont=bf,indention=2cm]{caption}

\begin{document}

test test test test test test test test test test test test test test
test test test test test test test test test test test test test test
test test test test test test test test test test test test test test
test test test test test test test test test test test test test test
test test test test test test test test test test test test test test
test test test test test test test test test test test test test test
test test test test test test test test test test test test test test
test test test test test test test test test test test test test test

\begin{figure}

\begin{center}
\framebox[5cm]{}
\end{center}

\caption{test test test test test test test test test test test test test
test test test test test test test test test test test test test test
test test test test test test test test test test test test test test}

\end{figure}

\end{document}
-------------------------------

Das Ergebnis habe ich als PDF angehängt.

Danke und Gruß
Uwe
\end{verbatim}

\begin{example}{labelfont=bf,indention=2cm}

\begin{figure}[H]
  \caption{test test test test test test test test test test test test test
  test test test test test test test test test test test test test test
  test test test test test test test test test test test test test test}
\end{figure}

\end{example}

%******************************************************************************%

\bug{05-03-23}{v3.0f}

\begin{verbatim}
Message-ID: <20050323213047.26829.qmail@web52507.mail.yahoo.com>
Date: Wed, 23 Mar 2005 13:30:46 -0800 (PST)
Subject: Fwd: [texhax] Problems specifying captions (with the captions package)
To: caption@sommerfeldt.net

Hello Mr. Sommerfeldt.
I am trying to use your LaTeX captions package and I'm
running into some trouble. I submitted a question to
the texhax list, but no one responded. I was wondering
if you would be so kind to address my question, which
I suspect for you would only take a minute. Thank you
very much and have a wonderful day.
Take Care,
Ken
P.S. Thanks for your work on LaTeX and for sharing
that work with others.

Note: forwarded message attached.

Date: Tue, 22 Mar 2005 21:37:29 -0800 (PST)
To: texhax@tug.org
Subject: [texhax] Problems specifying captions (with the captions package)

Hello all.

I'm trying to get table captions that are centered,
with the table identifier (i.e., TABLE 1) also in
caps, centered, and with no period after the
identifier. I'm trying to use the caption package to
make this happen, but I'm running into trouble. Most
of the options I need are not there and some of the
ones that are are not working out too well. I can
specify centered, but the table identifier part (e.g.,
TABLE 1 isn't actually centered). Furthermore, I
cannot figure out how to get the identifier in caps
(TABLE 1 rather than Table 1) . I've found a way to
cheat so that the caption itself is double spaced, but
this isn't convenient (by adding \newline commands
in). Any help would be greatly appreciate (it isn't my
idea to center and double space the table headings,
rather they are ridged guidelines I must follow).
Below is an example of what I'm coming up with (simply
past into the text editor and run LaTeX on it). I
can't find anything that discusses these problems on
line anywhere.

Thanks very much for any help.


% This is example code of what I'm getting.
\documentclass{amsart}
\usepackage{caption}
\captionsetup[table]{justification=centering,
labelsep=newline,font=small}
\begin{document}
\setcounter{table}{0}
\renewcommand{\thetable}{C-\arabic{table}}
\begin{table}
\begin{center}
\caption{\newline TABLE IDENTIFYING THE NUMBER OF
FAILURES THAT OCCURRED
IN THE\newline\newline SIMULATION STUDY FOR EACH OF
THE SITUATIONS
(CONDITION BY\newline\newline
SCENARIO)}\label{Failure.Table}
\begin{tabular}{|c|c|c|c|}
  \hline
  Condition & Scenario 1 & Scenario 2 & Scenario 3 \\
  \hline
  1 & 0 & 0 & 2 \\
  2 & 2 & 1 & 2 \\
  3 & 2 & 1 & 2 \\
  4 & 2 & 0 & 2 \\
  5 & 0 & 0 & 0 \\
  6 & 0 & 0 & 0 \\
  7 & 0 & 0 & 2 \\
  8 & 2 & 1 & 3 \\\hline
  9 & 0 & 1 & 2 \\
  10 & 2 & 1 & 0 \\
  11 & 0 & 1 & 3 \\
  12 & 1 & 2 & 4\\
  13 & 1 & 0 & 3 \\
  14 & 0 & 0 & 2 \\
  15 & 8 & 4 & 6 \\
  16 & 1 & 4 & 7 \\\hline
  17 & 0 & 0 & 0 \\
  18 & 0 & 0 & 0 \\
  19 & 0 & 1 & 6 \\
  20 & 1 & 1 & 0 \\
  21 & 0 & 0 & 3 \\
  22 & 0 & 0 & 1 \\
  23 & 1 & 4 & 16 \\
  24 & 2 & 3 & 12 \\
  \hline
\end{tabular}
\end{center}
\end{table}
\clearpage \pagebreak
\end{document}
% ----------------------------------
_______________________________________________
TeX FAQ: http://www.tex.ac.uk/faq
TeX newsgroup: http://groups.google.com/groups?group=comp.text.tex
Mailing list archives: http://tug.org/pipermail/texhax/
More links: http://tug.org/begin.html

Automated subscription management: http://tug.org/mailman/listinfo/texhax
Human mailing list managers: postmaster@tug.org
\end{verbatim}

\begin{example}{justification=centering,labelsep=newline,font=small,position=t}
\setcounter{table}{0}
\renewcommand{\thetable}{C-\arabic{table}}
\begin{table}[H]
\caption{TABLE IDENTIFYING THE NUMBER OF
FAILURES THAT OCCURRED
IN THE SIMULATION STUDY FOR EACH OF
THE SITUATIONS
(CONDITION BY SCENARIO)}
\centering
\begin{tabular}{|c|c|c|c|}
  \hline
  Condition & Scenario 1 & Scenario 2 & Scenario 3 \\
  \hline
  1 & 0 & 0 & 2 \\
  2 & 2 & 1 & 2 \\
  3 & 2 & 1 & 2 \\
  \hline
\end{tabular}
\end{table}
\end{example}

%******************************************************************************%

\bug{05-06-28}[v3.0f]{v3.0g}

\begin{verbatim}
Newsgroups: de.comp.text.tex
Subject: Re: Problem mit captcont
Date: Mon, 27 Jun 2005 14:45:31 +0200
Message-ID: <87psu8x91w.fsf@ID-223132.user.uni-berlin.de>

Axel Sommerfeldt <sommerfee@despammed.com> writes:

> Wenn irgendetwas nicht funktioniert oder du weitere Wünsche hast, zögere 
> bitte nicht, diese auch auszusprechen.

Ein Problem ist mir aufgefallen (Version 2005/06/22 v3.0f-BETA). Wenn
ich mit \include und \includeonly arbeite, so kreige ich die
Fehlermeldung

ERROR: LaTeX Error: No counter 'CF@figure' defined.

wenn ich den Abschnitt nicht übersetzte, in dem ich \ContinuedFloat
benutze. Wobei ich auch etwas überrascht war, das bei drei eingebundenen
Dateien foo.tex, bar.tex und baz.tax in foo.aux und bar.aux CF@figure
auf 2 gesetzt wurde, während es in baz.aux auf 0 gestzt wurde. Dabei
habe ich in foo.tex 2x \ContinuedFloat verwendet.

cheerio
ralf

Newsgroups: de.comp.text.tex
Subject: Re: Problem mit captcont
Date: Mon, 27 Jun 2005 22:31:22 +0200
Message-ID: <87k6kfy21x.fsf@ID-223132.user.uni-berlin.de>

Axel Sommerfeldt <sommerfee@despammed.com> writes:

>> ERROR: LaTeX Error: No counter 'CF@figure' defined.
>
> Auch, wenn du die aux-Dateien vorher löschst?

Dann nicht. Allerdings ist aus meiner Sicht das Vorhandensein und Lesen
der aux Dateien der eigentlich Witz bei include(only), da dadurch
Referenzen weiterhin möglich sind. Ich habe mal dazu ein Beispiel
gebastelt: 

%%%%%%%%%% main.tex %%%%%%%
\documentclass{article}
\usepackage{caption}[2005/06/22]
\usepackage{hyperref}
%\includeonly{bar}
\begin{document}
\include{foo}
\include{bar}
\include{baz}
\end{document}
%%%%%%%%%%%%%%%%%%%%%%%%%%%

%%%%%%%%%% foo.tex %%%%%%%%
\begin{figure}
  foo1
  \caption{foo1}\label{fig:foo1}
\end{figure}

\begin{figure}\ContinuedFloat
  foo2
  \caption{foo2}\label{fig:foo2}
\end{figure}
%%%%%%%%%%%%%%%%%%%%%%%%%%%

%%%%%%%%%% bar.tex %%%%%%%%
See figure~\ref{fig:foo1}.
%%%%%%%%%%%%%%%%%%%%%%%%%%%

%%%%%%%%%% baz.tex %%%%%%%%
\begin{figure}
  baz
  \caption{baz}\label{fig:baz}
\end{figure}
%%%%%%%%%%%%%%%%%%%%%%%%%%%

Dabei dient baz.tex nur zur Illustration des unterscheidlich gesetzten
Zählers in den aux Dateien dient. Bei einer ältern caption Version kann
ich nach einem Komplettdurchlauf auf \includeonly umsteigen, wobei der
Querverweis in bar.tex weiterhin funktioniert. Das geht aber nur, wenn
die aus Dateien aus dem Komplettdurchlauf vorhanden sind.

cheerio
ralf
\end{verbatim}

%******************************************************************************%

\bug{05-08-12}{v3.0h}

\begin{verbatim}
Newsgroups: de.comp.text.tex
Subject: Problem mit captionof und newfloat
Date: 12 Aug 2005 02:41:32 -0700
Message-ID: <1123839692.460760.151930@g43g2000cwa.googlegroups.com>

Hallo!

Ich möchte eine figure und den Inhalt einer mit \newfloat definierten
Umgebung zusammen halten, d.h. sie sollen gemeinsam in einer
Fließumgebung sitzen, aber ihre eigenen captions bekommen. Bei der
Suche in dieser Gruppe wurde ein ähnliches Problem schonmal
diskutiert, daraus habe ich die \captionof Definition von Axel
Sommerfeldt. Bei Nutzung von caption2 gehen die Formatierungen des
KOMA-Skripts verloren.

Die Probleme:
- der Abstand der caption zur Umgebung ist zu klein
- hyperlinks auf die eigene Umgebung (im Bsp. Algorithmus 1.2)
funktionieren nicht

Ich hoffe, jemand kann helfen,
Gruß,
Florian

hier das Beispiel:

\documentclass[a5paper,10pt,pdftex,DIVcalc,ngerman,headsepline,liststotoc]{scrbook}

\usepackage[latin1]{inputenc}
\usepackage[T1]{fontenc}
\usepackage{babel}

%\usepackage{caption2} % für captionof
% -> ändert KOMA-Formatierungen der captions

\setkomafont{caption}{\slshape}
\setkomafont{captionlabel}{\normalfont\bfseries}
\setcapwidth{.85\textwidth}

\usepackage{float}

\usepackage[linktocpage,bookmarksnumbered,bookmarksopen]{hyperref}

\floatstyle{komaabove}
\newfloat{algorithmus}{htb}{loa}[chapter]
\floatname{algorithmus}{Algorithmus}

\makeatletter
  \newcommand*\captionof[1]{\def\@captype{#1}%
    \@ifundefined{fst@#1}{}{%
      \@nameuse{fst@#1}%
      \@ifundefined{@float@setevery}{}{\@float@setevery{#1}}%
      \let\caption@fs@capt\@fs@capt
      \renewcommand\@fs@capt[2]{\egroup
        \vskip\abovecaptionskip
        \normalsize\caption@fs@capt{##1}{##2}%
        \vskip\belowcaptionskip
        \bgroup}}%
    \caption}
\makeatother

\listfiles
%%%%%%%%%%%%%%%%%%%%%%%%%%%%%%%%%%%%%%%%%%%%%%%%%%%%%%%%%%%
\begin{document}

\chapter{Bla}
bla bla bla bla bla bla bla bla bla bla bla bla bla bla bla bla bla bla
bla bla bla

\section{Abschnitt eins}

Siehe Algorithmus \ref{alg:eins}

\section{Abschnitt zwei}
Siehe Algorithmus \ref{alg:zwei} auf Seite \pageref{alg:zwei} und Abb.
\ref{fig:eins}.
\newpage

\begin{algorithmus}
\hrule
\smallskip
\centering
\begin{minipage}{.9\textwidth} % etwas schmaler als Fließtext
\begin{itemize}
	\item bla1
	\item bla2
	\item bla3
\end{itemize}
\end{minipage}
\smallskip
\hrule
\caption{Ein einzelner Code}
\label{alg:eins}
\end{algorithmus}

\begin{figure}
	\centering
\fbox{DUMMY}\par
\caption{Einzelnes Bild}
\label{fig:eins}
\end{figure}

\newpage

\begin{figure}
	\centering
\fbox{DUMMY}\par
\caption{Bild zusammen mit Code}
\label{fig:zwei}
\vspace{\floatsep} % zusaetzlicher Abstand zwischen zwei `floats'
\captionof{algorithmus}{Code zusammen mit Bild}
\label{alg:zwei}
\hrule
\smallskip
\centering
\begin{minipage}{.9\textwidth} % etwas schmaler als Fließtext
\begin{itemize}
	\item bla1
	\item bla2
	\item bla3
\end{itemize}
\end{minipage}
\smallskip
\hrule
\end{figure}

\end{document}
\end{verbatim}
%
% *File List*
% scrbook.cls    2004/09/16 v2.9t LaTeX2e KOMA document class
%scrlfile.sty    2004/09/16 v2.9t LaTeX2e KOMA package
%    bk10.clo    2004/02/16 v1.4f Standard LaTeX file (size option)
%typearea.sty    2004/09/16 v2.9t LaTeX2e KOMA package
%inputenc.sty    2004/02/05 v1.0d Input encoding file
%  latin1.def    2004/02/05 v1.0d Input encoding file
% fontenc.sty
%   t1enc.def    2004/02/22 v1.99f Standard LaTeX file
%   babel.sty    2004/02/19 v3.8a The Babel package
%ngermanb.ldf    2004/02/20 v2.6m new German support from the babel system
%   float.sty    2001/11/08 v1.3d Float enhancements (AL)
%hyperref.sty    2003/11/30 v6.74m Hypertext links for LaTeX
%  keyval.sty    1999/03/16 v1.13 key=value parser (DPC)
%  pd1enc.def    2003/11/30 v6.74m Hyperref: PDFDocEncoding definition (HO)
%hyperref.cfg    2003/03/08 v1.0 MiKTeX 'hyperref' configuration
%     url.sty    2004/03/15  ver 3.1  Verb mode for urls, etc.
% hpdftex.def    2003/11/30 v6.74m Hyperref driver for pdfTeX
%  pifont.sty    2004/09/15 PSNFSS-v9.2 Pi font support (SPQR)
%    upzd.fd    2001/06/04 font definitions for U/pzd.
%    upsy.fd    2001/06/04 font definitions for U/psy.
% nameref.sty    2003/12/03 v2.21 Cross-referencing by name of section float_captionof.out
%float_captionof.out
%  t1cmss.fd    1999/05/25 v2.5h Standard LaTeX font definitions
%  omscmr.fd    1999/05/25 v2.5h Standard LaTeX font definitions
%\end{verbatim}

%******************************************************************************%

\bug{05-11-30}{v3.0i}

\begin{verbatim}
Message-ID: <20051201044932.25609.qmail@web54204.mail.yahoo.com>
Date: Wed, 30 Nov 2005 20:49:32 -0800 (PST)
Subject: caption problem
To: caption@sommerfee.de

Axel,

I hope all is well.

I am very close to releasing my updated
epslatex document, and I noticed an 
interaction with the rotating and 
caption packages.

Attached is a pdflatex file that uses a
\rotcaption command.  When the caption
package is not used, the caption is 
typeset with a short width (so after
rotation it takes up very little 
vertical space).  When I use the caption
package, the caption has a very large 
width (so after rotation it takes up a
lot of vertical space).

Look through the attached file (with and
without \usepackage{caption}) and let
me know if you have any questions.

Thanks,
Keith


\documentclass[letterpaper,11pt,oneside,pdftex]{article}
\usepackage{rotating}
\usepackage{graphicx}
% \usepackage{caption}

\begin{document}

\noindent
The following commands
\begin {verbatim}
     \begin{figure}[!ht]
       \centering
       \begin{minipage}[c]{3cm}
          \hfill\includegraphics[width=4cm,angle=90]{graphic}
       \end{minipage}%
       \hspace{20pt}%
       \begin{minipage}[c]{30pt}
          % \captionsetup{width=4cm}
          \rotcaption{Rotcaption Caption}
          \label{fig:rotcaption}
       \end{minipage}
     \end{figure}
\end {verbatim}
produce Figure~\ref{fig:rotcaption}.

     \begin{figure}[!ht]
       \centering
       \begin{minipage}[c]{3cm}
          \hfill\includegraphics[width=4cm,angle=90]{graphic}
       \end{minipage}%
       \hspace{20pt}%
       \begin{minipage}[c]{30pt}
          % \captionsetup{width=4cm}
          \rotcaption{Rotcaption Caption}
          \label{fig:rotcaption}
       \end{minipage}
     \end{figure}

Text text text text text text text text text text text text text text.
Text text text text text text text text text text text text text text.
Text text text text text text text text text text text text text text.
Text text text text text text text text text text text text text text.
Text text text text text text text text text text text text text text.
Text text text text text text text text text text text text text text.

Text text text text text text text text text text text text text text.
Text text text text text text text text text text text text text text.
Text text text text text text text text text text text text text text.
Text text text text text text text text text text text text text text.
Text text text text text text text text text text text text text text.
Text text text text text text text text text text text text text text.
Text text text text text text text text text text text text text text.

Text text text text text text text text text text text text text text.
Text text text text text text text text text text text text text text.
Text text text text text text text text text text text text text text.
Text text text text text text text text text text text text text text.
Text text text text text text text text text text text text text text.
Text text text text text text text text text text text text text text.
Text text text text text text text text text text text text text text.
Text text text text text text text text text text text text text text.
Text text text text text text text text text text text text text text.
Text text text text text text text text text text text text text text.
Text text text text text text text text text text text text text text.

Text text text text text text text text text text text text text text.
Text text text text text text text text text text text text text text.
Text text text text text text text text text text text text text text.
Text text text text text text text text text text text text text text.
Text text text text text text text text text text text text text text.
Text text text text text text text text text text text text text text.

Text text text text text text text text text text text text text text.
Text text text text text text text text text text text text text text.
Text text text text text text text text text text text text text text.
Text text text text text text text text text text text text text text.
Text text text text text text text text text text text text text text.
Text text text text text text text text text text text text text text.

Text text text text text text text text text text text text text text.
Text text text text text text text text text text text text text text.
Text text text text text text text text text text text text text text.
Text text text text text text text text text text text text text text.
Text text text text text text text text text text text text text text.
Text text text text text text text text text text text text text text.

\end{document}
\end{verbatim}

%******************************************************************************%

\bug{06-01-03}{v3.0i}

\begin{verbatim}
Date: Tue, 03 Jan 2006 09:47:44 +0100
Newsgroups: de.comp.text.tex
Subject: KOMA-Script (scrbook) und subfig
Message-ID: <43ba3a30$1@news1.ethz.ch>

Erst mal ein Frohes Neues allerseits!

Bei mir tritt mal wieder ein Problem mit scrbook und subfig auf, mit dem
Minimalbeispiel unten bekomme ich folgende Fehlermeldung (latex und
pdflatex unter Miktex 2.4 (aktuell)):

! Package keyval Error: No value specified for parskip.

See the keyval package documentation for explanation.
Type  H <return>  for immediate help.
 ...

l.9   \subfloat
               {a}
?


An dem bekannten Problem mit dem caption-Paket scheint es nicht direkt
zu liegen, da \usepackage[caption=false]{subfig} zum gleichen Fehler führt.

Auf meinem Mac unter tetex3 läuft das ganze dann interessanterweise wieder.

Minimalbeispiel:

\documentclass[parskip]{scrbook}
\usepackage{subfig}
\listfiles

\begin{document}

\begin{figure}
  \subfloat{a}
  \qquad
  \subfloat{b}
\end{figure}

\end{document}

das pdflatex-logfile dazu habe ich mal angehängt.
\end{verbatim}

%\begin{verbatim}
%This is pdfeTeX, Version 3.141592-1.21a-2.2 (MiKTeX 2.4) (preloaded format=latex 2006.1.3)  3 JAN 2006 09:47
%entering extended mode
%**test.tex
%(test.tex
%LaTeX2e <2003/12/01>
%Babel <v3.8g> and hyphenation patterns for ukenglish, english, german, ngerman,
% dumylang, nohyphenation, loaded.
%(C:\program files\texmf\tex\latex\koma-script\scrbook.cls
%Document Class: scrbook 2004/09/16 v2.9t LaTeX2e KOMA document class
%(C:\program files\texmf\tex\latex\koma-script\scrlfile.sty
%Package: scrlfile 2004/09/16 v2.9t LaTeX2e KOMA package
%
%Package scrlfile, 2004/09/16 v2.9t LaTeX2e KOMA package
%                  Copyright (C) Markus Kohm
%
%) (C:\program files\texmf\tex\latex\base\bk11.clo
%File: bk11.clo 2004/02/16 v1.4f Standard LaTeX file (size option)
%)
%(C:\program files\texmf\tex\latex\koma-script\typearea.sty
%Package: typearea 2004/09/16 v2.9t LaTeX2e KOMA package
%
%Package typearea, 2004/09/16 v2.9t LaTeX2e KOMA package
%                  Copyright (C) Frank Neukam, 1992-1994
%                  Copyright (C) Markus Kohm, 1994-2002
%
%\ta@bcor=\skip41
%\ta@div=\count79
%\ta@hblk=\skip42
%\ta@vblk=\skip43
%\ta@temp=\skip44
%Package typearea Info: These are the values describing the layout:
%(typearea)             DIV  = 10
%(typearea)             BCOR = 0.0pt
%(typearea)             \paperwidth      = 597.50793pt
%(typearea)              \textwidth      = 418.25555pt
%(typearea)              \columnwidth    = 0.0pt
%(typearea)              \columnsep      = 0.0pt
%(typearea)              DIV-departure   = -6/100
%(typearea)              \evensidemargin = 47.2316pt
%(typearea)              \oddsidemargin  = -12.5192pt
%(typearea)             \paperheight     = 845.04694pt
%(typearea)              \textheight     = 595.80026pt
%(typearea)              \topmargin      = -25.16531pt
%(typearea)              \headheight     = 17.0pt
%(typearea)              \headsep        = 20.40001pt
%(typearea)              \topskip        = 11.0pt
%(typearea)              \footskip       = 47.60002pt
%(typearea)              \baselineskip   = 13.6pt
%(typearea)              on input line 633.
%)
%\c@part=\count80
%\c@chapter=\count81
%\c@section=\count82
%\c@subsection=\count83
%\c@subsubsection=\count84
%\c@paragraph=\count85
%\c@subparagraph=\count86
%\c@figure=\count87
%\c@table=\count88
%\abovecaptionskip=\skip45
%\belowcaptionskip=\skip46
%\c@pti@nb@sid@b@x=\box26
%\bibindent=\dimen102
%) (C:\program files\texmf\tex\latex\subfig\subfig.sty
%Package: subfig 2005/06/28 ver: 1.3 subfig package
%
%(C:\program files\texmf\tex\latex\graphics\keyval.sty
%Package: keyval 1999/03/16 v1.13 key=value parser (DPC)
%\KV@toks@=\toks14
%)
%(C:\program files\texmf\tex\latex\caption\caption.sty
%Package: caption 2005/10/24 v3.0h Customising captions (AR)
%
%(C:\program files\texmf\tex\latex\caption\caption3.sty
%Package: caption3 2005/10/24 v3.0h caption3 kernel (AR)
%\captionmargin=\dimen103
%\captionmarginx=\dimen104
%\captionwidth=\dimen105
%\captionindent=\dimen106
%\captionparindent=\dimen107
%\captionhangindent=\dimen108
%)
%Package caption Info: Option `parskip' ignored on input line 146.
%)
%\c@KVtest=\count89
%\sf@farskip=\skip47
%\sf@captopadj=\dimen109
%\sf@capskip=\skip48
%\sf@nearskip=\skip49
%\c@subfigure=\count90
%\c@subfigure@save=\count91
%\c@lofdepth=\count92
%\c@subtable=\count93
%\c@subtable@save=\count94
%\c@lotdepth=\count95
%\sf@top=\skip50
%\sf@bottom=\skip51
%) (test.aux)
%LaTeX Font Info:    Checking defaults for OML/cmm/m/it on input line 5.
%LaTeX Font Info:    ... okay on input line 5.
%LaTeX Font Info:    Checking defaults for T1/cmr/m/n on input line 5.
%LaTeX Font Info:    ... okay on input line 5.
%LaTeX Font Info:    Checking defaults for OT1/cmr/m/n on input line 5.
%LaTeX Font Info:    ... okay on input line 5.
%LaTeX Font Info:    Checking defaults for OMS/cmsy/m/n on input line 5.
%LaTeX Font Info:    ... okay on input line 5.
%LaTeX Font Info:    Checking defaults for OMX/cmex/m/n on input line 5.
%LaTeX Font Info:    ... okay on input line 5.
%LaTeX Font Info:    Checking defaults for U/cmr/m/n on input line 5.
%LaTeX Font Info:    ... okay on input line 5.
%
%(C:\program files\texmf\tex\latex\ms\ragged2e.sty
%Package: ragged2e 2003/03/25 v2.04 ragged2e Package (MS)
%
%(C:\program files\texmf\tex\latex\ms\everysel.sty
%Package: everysel 1999/06/08 v1.03 EverySelectfont Package (MS)
%LaTeX Info: Redefining \selectfont on input line 125.
%)
%\CenteringLeftskip=\skip52
%\RaggedLeftLeftskip=\skip53
%\RaggedRightLeftskip=\skip54
%\CenteringRightskip=\skip55
%\RaggedLeftRightskip=\skip56
%\RaggedRightRightskip=\skip57
%\CenteringParfillskip=\skip58
%\RaggedLeftParfillskip=\skip59
%\RaggedRightParfillskip=\skip60
%\JustifyingParfillskip=\skip61
%\CenteringParindent=\skip62
%\RaggedLeftParindent=\skip63
%\RaggedRightParindent=\skip64
%\JustifyingParindent=\skip65
%)
%Package caption Info: KOMA-Script class detected on input line 5.
%
%
%! Package keyval Error: No value specified for parskip.
%
%See the keyval package documentation for explanation.
%Type  H <return>  for immediate help.
% ...                                              
%                                                  
%l.8   \subfloat
%               {a}
%? 
%
%! Package keyval Error: No value specified for parskip.
%
%See the keyval package documentation for explanation.
%Type  H <return>  for immediate help.
% ...                                              
%                                                  
%l.10   \subfloat
%                {b}
%? 
%[1{psfonts.map}] (test.aux)
%
% *File List*
% scrbook.cls    2004/09/16 v2.9t LaTeX2e KOMA document class
%scrlfile.sty    2004/09/16 v2.9t LaTeX2e KOMA package
%    bk11.clo    2004/02/16 v1.4f Standard LaTeX file (size option)
%typearea.sty    2004/09/16 v2.9t LaTeX2e KOMA package
%  subfig.sty    2005/06/28 ver: 1.3 subfig package
%  keyval.sty    1999/03/16 v1.13 key=value parser (DPC)
% caption.sty    2005/10/24 v3.0h Customising captions (AR)
%caption3.sty    2005/10/24 v3.0h caption3 kernel (AR)
%ragged2e.sty    2003/03/25 v2.04 ragged2e Package (MS)
%everysel.sty    1999/06/08 v1.03 EverySelectfont Package (MS)
% ***********
%
% ) 
%Here is how much of TeX's memory you used:
% 1380 strings out of 95509
% 20066 string characters out of 1189015
% 77790 words of memory out of 1081545
% 4495 multiletter control sequences out of 60000
% 4293 words of font info for 16 fonts, out of 1000000 for 2000
% 22 hyphenation exceptions out of 8191
% 43i,7n,38p,175b,162s stack positions out of 5000i,500n,10000p,200000b,32768s
%PDF statistics:
% 7 PDF objects out of 300000
% 0 named destinations out of 300000
% 1 words of extra memory for PDF output out of 65536
%<C:\program files\texmf\fonts\type1\bluesky\cm\cmr10.pfb>
%Output written on test.pdf (1 page, 4252 bytes).
%\end{verbatim}

%******************************************************************************%

\bug{06-01-08}{v3.0i}

\begin{verbatim}
Message-ID: <20060108231644.29592.qmail@web54208.mail.yahoo.com>
Date: Sun, 8 Jan 2006 15:16:44 -0800 (PST)
Subject: Re: Some remarks regarding the caption package within your  epslatex document
To: Axel Sommerfeldt <axel.sommerfeldt@arcor.de>
In-Reply-To: <5.2.1.1.2.20051229171835.01a04ad0@pop3.arcor.de>

Axel,

Happy New Year...

Thank you very much for those comments.  I have incorporated 
them in my document and will be releasing a typo-fixed version 
in the next few days.

I have also attached a latex source file "test_cap.tex" and its 
output "test_cap.pdf" which shows a couple questions I have 
with the caption package.

The first example shows a custom caption format where I don't
understand why I get a newline before the \gg command.

The second example shows a figure that extends into the
margin.  The caption is centered in the document's text width, 
not in the wider figure width.  I am not sure if you consider
this a bug or not.

Thanks,
Keith


\documentclass{article}
\usepackage{graphicx}
\usepackage{caption}
\usepackage{verbatim}
\usepackage{geometry}

\geometry{left=0.8in, right=0.9in, top=0.6in, bottom=0.6in,
          includeall, reversemp, %
          marginparwidth=2in, marginparsep=0.15in,  % marginal notes
          headheight=0pt,       % default 12pt
          headsep=0pt,          % default 20pt
          footskip=24pt}        % default 36pt

%-----------------------------------------------------------------
% \begin{narrow}{1.0in}{0.5in}   produces text which is narrowed
% by 1.0 on left margin and 0.5 inches on right margin
% \begin{narrow}{-1.0in}{-0.5in} produces text which is widened
% by 1.0 on left margin and 0.5 inches on right margin
% Narrow environments can be nested and are ended by \end{narrow}
%-----------------------------------------------------------------
\newenvironment{narrow}[2]{%
 \begin{list}{}{%
  \setlength{\topsep}{0pt}%
  \setlength{\leftmargin}{#1}%
  \setlength{\rightmargin}{#2}%
  \setlength{\listparindent}{\parindent}%
  \setlength{\itemindent}{\parindent}%
  \setlength{\parsep}{\parskip}%
 }%
\item[]}{\end{list}}

\DeclareCaptionFormat{angle}{\ensuremath{\ll}#1#2#3\ensuremath{\gg}}

\begin{document}

\newlength{\marginwidth}
\setlength{\marginwidth}{\marginparwidth}
\addtolength{\marginwidth}{\marginparsep}

\section{Custom Caption Format}

For example, the following command in the document's preamble
\begin {verbatim}
    \DeclareCaptionFormat{angle}{\ensuremath{\ll}#1#2#3\ensuremath{\gg}}
\end {verbatim}
defines the \texttt{angle} format with double angle brackets 
\verb+\ll+ and \verb+\gg+ surrounding the caption.
The following code
\begin {verbatim}
   \captionsetup{format=angle}
   \caption{This Caption has a Custom Format}
\end {verbatim}
produces the caption in Figure~\ref{fig:caption:angle}.
\begin{figure}[htbp]
    \centering
    \captionsetup{format=angle}
    \includegraphics[width=2in]{wide}
    \captionof{figure}{This Caption has a Custom Format}
    \label{fig:caption:angle}
\end{figure}

\section{Wide Figures}
When the narrow environment is defined as
\begin {verbatim}
     \newenvironment{narrow}[2]{% 
      \begin{list}{}{% 
       \setlength{\topsep}{0pt}% 
       \setlength{\leftmargin}{#1}% 
       \setlength{\rightmargin}{#2}% 
       \setlength{\listparindent}{\parindent}% 
       \setlength{\itemindent}{\parindent}% 
       \setlength{\parsep}{\parskip}% 
      }% 
     \item[]}{\end{list}} 
\end {verbatim}
The following code 
\begin {verbatim}
     \begin{figure}
     \begin{narrow}{-2in}{0in}
        \includegraphics[width=\linewidth]{wide}
        \caption{This is a wide figure}
     \end{narrow}
     \end{figure}
\end {verbatim}
uses this \texttt{narrow} environment to 
make the figure extend 2~inches into the left margin,
producing Figure~\ref{fig:wide}.

\begin{figure}[htbp]
\begin{narrow}{-2in}{0in}
    \includegraphics[width=\linewidth]{wide}
    \caption{This is a wide figure}
    \label{fig:wide}
\end{narrow}
\end{figure}

\end{document}
\end{verbatim}

%******************************************************************************%

\bug{06-02-22}{v3.0j}

\begin{verbatim}
Date:         Wed, 22 Feb 2006 16:03:30 +0100
Reply-To:     German Language TeX Users Group Communication List
              <TEX-D-L@LISTSERV.DFN.DE>
Sender:       German Language TeX Users Group Communication List
              <TEX-D-L@LISTSERV.DFN.DE>
Subject:      caption und hyperref

Moin,

ein Schuß ins Blaue, bevor ich mir ein Minimalbeispiel
aus den Rippen schneide. Vielleicht hat jemand eine 
frisch polierte Kristallkugel dabei ...

In welcher Richtung habe ich zu suchen, wenn die unten
angehängte Fehlermeldung beim LaTeXen auftaucht? Der
Hinweis, in der caption-Doku zu suchen, hat mich leider
nicht erleuchtet.

Gruß und vielen Dank im voraus

Jürgen



(c:/Program Files/TeXLive/texmf-local/tex/latex/subfig/subfig.sty
Package: subfig 2004/01/28 ver: 1.2 subfig package

(c:/Program Files/TeXLive/texmf-local/tex/latex/caption/caption.sty
Package: caption 2006/01/12 v3.0i Customising captions (AR)

(c:/Program Files/TeXLive/texmf-local/tex/latex/caption/caption3.sty
Package: caption3 2006/01/12 v3.0i caption3 kernel (AR)
\captionmargin=\dimen163
\captionmarginx=\dimen164
\captionwidth=\dimen165
\captionindent=\dimen166
\captionparindent=\dimen167
\captionhangindent=\dimen168
)

! Package caption Error: No value specified for hyperref.

See the caption package documentation for explanation.
Type  H <return>  for immediate help.
 ...

l.122 \ProcessOptionsWithKV{caption}

?
Package caption Info: float package v1.3 (or newer) detected on input line 263.



--
FAQ: http://www.dante.de/faq/de-tex-faq/

Durchsuchbares Archiv: http://tinyurl.com/cdcb6
Unsubscribe/Verwaltung: http://tinyurl.com/b9tod
\end{verbatim}

%******************************************************************************%

\bug{06-03-08}{v3.1i}

\begin{verbatim}
Newsgroups: de.comp.text.tex
Subject: Bilder nebeneinander und Abstände von \captionof
Date: Wed, 08 Mar 2006 16:23:35 +0100
Message-ID: <dumstu$gg8$1@sagnix.uni-muenster.de>

Hallo,

ich möchte zwei Bilder nebeneinander setzten, oben bündig ausrichten und
mit je einer Unterschrift versehen.
Prinzipiell klappt das auch, bis darauf, dass der Abstand zwischen Bild
und Unterschrift nicht immer gleich groß ist, sondern davon abhängt, ob
die Unterschrift ein- oder mehrzeilig ist.

Wie kann man das beheben?

Minimalbsp:

\documentclass{scrartcl}
\usepackage{caption}

\begin{document}
\noindent\begin{minipage}{\linewidth}
  \parbox[t]{0.3\linewidth}{%
    \vspace{0pt}
    \centering
    \rule{3cm}{3cm}
    \captionof{figure}{short caption}
    }%
  \parbox[t]{0.3\linewidth}{%
    \vspace{0pt}
    \centering
    \rule{3cm}{3cm}
    \captionof{figure}{long caption long caption long caption long
                       caption long caption}
  }%
  \parbox[t]{0.3\linewidth}{%
    \vspace{0pt}
    \centering
    \rule{3cm}{3cm}
    \captionof{figure}{longer caption longer caption longer caption
                       longer caption longer caption}
  }%
\end{minipage}
\end{document}

Ein \listfiles liefert mir

 *File List*
scrartcl.cls    2003/01/31 v2.9n LaTeX2e KOMA document class
scrlfile.sty    2003/01/31 v2.9n LaTeX2e KOMA package
  size11.clo    2001/04/21 v1.4e Standard LaTeX file (size option)
typearea.sty    2003/01/31 v2.9n LaTeX2e KOMA package
 caption.sty    2006/01/12 v3.0i Customising captions (AR)
caption3.sty    2006/01/12 v3.0i caption3 kernel (AR)
  keyval.sty    1999/03/16 v1.13 key=value parser (DPC)
ragged2e.sty    2003/02/24 v2.02 ragged2e Package (MS)
everysel.sty    1999/06/08 v1.03 EverySelectfont Package (MS)
 ***********


Besten Dank
Christoph
\end{verbatim}

%******************************************************************************%

\bug{06-03-14}{v3.1}

\begin{verbatim}
Newsgroups: de.comp.text.tex
Subject: \hrule vs. addmargin-Umgebung
Date: 14 Mar 2006 15:05:17 GMT
Message-ID: <47o4ddFgdt0iU1@news.dfncis.de>

Guten Tag,

ich habe hier ein eigenes \caption Layout, das die jeweilige
Beschriftung zwischen zwei \hrule setzt. Diese \caption benutze
ich jetzt neuerdings auch in addmargin Umgebungen, da meine
Tabellen den Satzspiegel verlassen.

Jetzt habe ich das Problem, dass allerdings die \hrule innerhalb
der addmargin-Umgebung nicht länger geworden ist, d.h. also zu kurz
ist.

Wie könnte ich a) mein \caption Layout umdefinieren, so dass ich
entscheiden kann, ob ich in einer addmargin-Umgebung bin oder nicht
(um dann ggfs. eine entsprechend längere \hrule zu setzen) oder b)
eine eigene \hrule basteln, die automatisch länger wird, falls wir
in einer addmargin Umgebung sind?

Herzlichen Gruß,

Michael.
\end{verbatim}

%******************************************************************************%

\bug{06-03-21}{v3.0j}

\begin{verbatim}
Newsgroups: de.comp.text.tex
Subject: [KOMA+caption]: Seitliche Tabellenüberschriften
Date: Mon, 20 Mar 2006 16:22:44 +0100
Message-ID: <uu09tp2d7.fsf@uwe-siart.de>

Hallo,

ich weiß, dass sich KOMA-Klassen und das 'caption'-Paket zwar schon sehr
gut, aber noch nicht in allen Punkten vertragen. Markus und Axel haben
mir vor etwa zwei Jahren auch erklärt, was alles getan wird, um das zu
ändern und dass zur Zeit der einzige Weg sein kann, entweder auf KOMA
oder auf 'caption' zu verzichten.

Weil ich in einem großen Dokument scrbook.cls und das 'subfig'-Paket
(und damit auch das 'caption'-Paket) verwende, möchte ich fragen, ob es
für folgenden Effekt nicht trotzdem einen Workaround gibt:

Wenn KOMA + 'caption' verwendet wird und man Tabellenüberschriften haben
möchte, dann kommt es bei seitlichen Überschriften mittels der
{captionbeside}-Umgebung zu einer Fehlpositionierung der Legende: Sie
wird zu hoch platziert.

Den Effekt sieht man hier, vergleiche mit/ohne {caption}:

% --------------------------------------------
\documentclass{scrartcl}
\usepackage[tableposition=top]{caption}
\begin{document}
\begin{table}
  \begin{captionbeside}{Tabellen\"uberschrift}
    \begin{tabular}{ll}\hline
      Zelle & Zelle \\ \hline
    \end{tabular}
  \end{captionbeside}
\end{table}
\end{document}
% --------------------------------------------

Gibt einen Trick, das zu retten?

 *File List*
scrartcl.cls    2004/09/16 v2.9t LaTeX2e KOMA document class
scrlfile.sty    2004/09/16 v2.9t LaTeX2e KOMA package
  size11.clo    2004/02/16 v1.4f Standard LaTeX file (size option)
typearea.sty    2004/09/16 v2.9t LaTeX2e KOMA package
 caption.sty    2006/01/12 v3.0i Customising captions (AR)
caption3.sty    2006/01/12 v3.0i caption3 kernel (AR)
  keyval.sty    1999/03/16 v1.13 key=value parser (DPC)
ragged2e.sty    2003/03/25 v2.04 ragged2e Package (MS)
everysel.sty    1999/06/08 v1.03 EverySelectfont Package (MS)
 ***********

-- 
Uwe
\end{verbatim}

%******************************************************************************%

\bug{06-05-15}{v3.0k}

\begin{verbatim}
Date: Mon, 15 May 2006 10:04:33 -0400
Newsgroups: comp.text.tex
Subject: wrapfig and caption
Message-ID: <44688a65$0$558$b45e6eb0@senator-bedfellow.mit.edu>

Hi,

I found that when I use both wrapfig, caption, and setspace packages 
together, the line spacing in caption changes to the spacing of the 
text.  I want to use different line spacing for my captions.  I found an 
instruction in caption package but it does not work with either wrapfig 
or setspace.

Thanks.

Dan
\end{verbatim}

%******************************************************************************%

\bug{06-05-31}{-- , Bug is in subfig package, not in caption package}

\begin{verbatim}
Newsgroups: comp.text.tex
Subject: new caption (3.0j) and adjustwidth (regression?)
Date: Wed, 31 May 2006 21:36:24 +0000 (UTC)
Message-ID: <e5l28o$na9$1@netlx020.civ.utwente.nl>

Hello,

I am trying set two floats next to each other with subfig. This works
all ok, unless it's wider than \textwidth. Below is a minimal example to
demonstrate. I am using adjustwidth of package chngpage.

Using caption version 3.0c it all looks ok (the main caption of figure 1.1
is not exactly centered though):

 *File List*
 article.cls    2004/02/16 v1.4f Standard LaTeX document class
  size10.clo    2004/02/16 v1.4f Standard LaTeX file (size option)
 caption.sty    2004/07/16 v3.0c Customising captions (AS)
  keyval.sty    1999/03/16 v1.13 key=value parser (DPC)
  subfig.sty    2005/06/28 ver: 1.3 subfig package
chngpage.sty    2003/08/10 v1.2 change page layout
ragged2e.sty    2003/03/25 v2.04 ragged2e Package (MS)
everysel.sty    1999/06/08 v1.03 EverySelectfont Package (MS)
 ***********

But using newer version 3.0j, I get the following output. The overfull
boxes are both subcaptions. The first exceeds in the second. The second
exceeds the margin.

Overfull \hbox (113.811pt too wide) in paragraph at lines 22--22
[] []

Overfull \hbox (113.811pt too wide) in paragraph at lines 27--27
[] []
[1{/var/lib/texmf/fonts/map/pdftex/updmap/pdftex.map}] (./subfl.aux)

 *File List*
 article.cls    2004/02/16 v1.4f Standard LaTeX document class
  size10.clo    2004/02/16 v1.4f Standard LaTeX file (size option)
 caption.sty    2006/03/21 v3.0j Customising captions (AR)
caption3.sty    2006/03/16 v3.0j caption3 kernel (AR)
  keyval.sty    1999/03/16 v1.13 key=value parser (DPC)
  subfig.sty    2005/06/28 ver: 1.3 subfig package
chngpage.sty    2003/08/10 v1.2 change page layout
ragged2e.sty    2003/03/25 v2.04 ragged2e Package (MS)
everysel.sty    1999/06/08 v1.03 EverySelectfont Package (MS)
 ***********

But the release notes of caption 3.0i says:
Now works correctly inside list-based environments like "narrow" or "wide"

The real reason I am trying this, is because I want to put two source
listings (lstlisting environment) inside subfloat. I created the
following package to achieve this:

http://composestar.cvs.sourceforge.net/composestar/home/thesis2006/style/lstsub.sty?view=markup

It creates a newsubfloat for lstlisting.

Can anyone shed some light on this problem?

Cheers
 -Olaf

The is the minimal example (which goes wrong with caption v3.0j):


\documentclass{article}

\usepackage{caption}
\usepackage{subfig}
\usepackage{chngpage}

\listfiles

\begin{document}

\section{Hello}

This is one paragraph of text.
Foo foo foo foo foo foo foo foo.

\begin{figure}[t]
  \begin{adjustwidth}{-2cm}{-2cm}
  \centering
  \subfloat[One subcaption. Let's make this subcaption a bit longer.]{%
    \label{fig:example_one_a}%
    \parbox{7cm}{One one one one one one one one one one one one one one one one.}%
  }%
  \qquad
  \subfloat[Another subcaption.]{%
    \label{fig:example_one_b}%
    \parbox{7cm}{Two two two two two two two two two two two two two two two two.}%
  }
  \caption{The main caption is here.}%
  \label{fig:example_one}%
  \end{adjustwidth}%
\end{figure}

This is also one paragraph of text to conclude this example, see Figure~\ref{fig:example_one}.
Bar bar bar bar bar bar bar bar.

\begin{figure}[b]
  \centering
  \subfloat[One subcaption. Let's make this subcaption a bit longer.]{%
    \label{fig:example_two_a}%
    \parbox{5cm}{One one one one one one one one one one one one one one one one.}%
  }%
  \qquad
  \subfloat[Another subcaption.]{%
    \label{fig:example_two_b}%
    \parbox{5cm}{Two two two two two two two two two two two two two two two two.}%
  }
  \caption{The main caption is here.}%
  \label{fig:example_two}%
\end{figure}

\end{document}
\end{verbatim}

\begin{example}{}

\begin{figure}[H]
  \begin{adjustwidth}{-2cm}{-2cm}
  \centering
  \subfloat[One subcaption. Let's make this subcaption a bit longer.]{%
    \label{fig:example_one_a}%
    \parbox{7cm}{One one one one one one one one one one one one one one one one.}%
  }%
  \qquad
  \subfloat[Another subcaption.]{%
    \label{fig:example_one_b}%
    \parbox{7cm}{Two two two two two two two two two two two two two two two two.}%
  }
  \caption{The main caption is here.}%
  \label{fig:example_one}%
  \end{adjustwidth}%
\end{figure}

\begin{figure}[H]
  \centering
  \subfloat[One subcaption. Let's make this subcaption a bit longer.]{%
    \label{fig:example_two_a}%
    \parbox{5cm}{One one one one one one one one one one one one one one one one.}%
  }%
  \qquad
  \subfloat[Another subcaption.]{%
    \label{fig:example_two_b}%
    \parbox{5cm}{Two two two two two two two two two two two two two two two two.}%
  }
  \caption{The main caption is here.}%
  \label{fig:example_two}%
\end{figure}

\end{example}

%******************************************************************************%

\bug{06-10-12}{v3.1}

\begin{verbatim}
Newsgroups: de.comp.text.tex
Subject: Unverträglichkeit von caption.sty und captionsbeside (KOMAScript)?
Date: 12 Oct 2006 06:17:13 -0700
Message-ID: <1160659033.196840.138750@c28g2000cwb.googlegroups.com>

Moin,

mein Problem mit captionbeside aus KOMAScript zeigt das
unten stehende Beispiel (sofern ich es nicht wieder vergesse):
In der enumerate-Umgebung fällt die Gesamt-Breite von
Abbildung und caption zu gering aus. Ich habe schon versucht,
in einer Gruppe direkt um die Umgebung die \linewidth wieder
auf \textwidth zu setzen, aber dieser naive Ansatz war
offensichtlich untauglich. Das Problem tritt allerdings erst auf,
wenn ich \usepackage{caption} im Header habe. Und auf
dieses Paket möchte ich nicht so gerne verzichten. Wie kann
ich captionbeside überreden, die ganze Textbreite zu nutzen?

Gruß

Jürgen

P. S.: Diese Frage ging gestern früh schon an TeX-D-L, bisher
allerdings ohne Erfolg. Da es klein wenig eilt, hier der zweite
Ansatz.


\documentclass[parskip]{scrartcl}
\usepackage{caption}
\usepackage{blindtext}
\begin{document}
\begin{enumerate}
\item \blindtext
%
\begin{figure}
\begin{captionbeside}{Diese caption steht neben der Abbildung und so
weiter und so fort}
 \rule{0.5\linewidth}{10\baselineskip}
\end{captionbeside}
\end{figure}
%
\blindtext
\end{enumerate}
\end{document}
\end{verbatim}

%******************************************************************************%

\bug{07-01-19}{v3.0l}

\begin{verbatim}
To: caption@sommerfee.de
Subject: Caption version 2006/03/16 v3.0j problem and solution
Date: 19 Jan 2007 12:39:22 +0100
Message-ID: <418xfzfdp1.fsf@obel19.ida.liu.se>

One of our PhD students encountered a problem with the caption package
when used in conjunction with the combine package.

When using caption and combine together, \caption causes LaTeX to go
into an endless loop.

Caption uses \AtBeginDocument to patch the \caption macro, but does
not protect against defining it twice. When using combine,
AtBeginDocument code is expanded once for the master document and once
for each included document. On the second expansion, \caption@old is
redefined as the patched \caption, so when \caption is expanded, a
never-ending sequence of \caption@caption is produced.

I suggest you don't patch \caption if \caption@old is defined, similar
to how \caption is not patched if \cc@caption is defined. Something
like this ought to get the job done (though this patch is *not* tested
since I patched the sty file directly).


--- caption.dtx.old     Fri Jan 19 12:36:41 2007
+++ caption.dtx Fri Jan 19 10:51:28 2007
@@ -3701,6 +3701,7 @@
 %
 %    \begin{macrocode}
   \@ifundefined{cc@caption}{%
+  \@ifundefined{caption@old}{%
 %    \end{macrocode}
 %
 % \begin{macro}{\caption}
@@ -3730,7 +3731,7 @@
 % \end{macro}
 %
 %    \begin{macrocode}
-  }{%
+  }}{%
 %    \end{macrocode}
 %
 % \changes{v3.0c}{2004/07/16}{Bugfix 04-07-15: \package{captcont} support fixed}



-- 
David Byers.
\end{verbatim}

%******************************************************************************%

\bug{07-01-20}{v3.0l}

\begin{verbatim}
Message-ID: <000f01c73c80$d3ca4f00$3e653ac2@home196ib52sok>
To: "Axel Sommerfeldt" <axel.sommerfeldt@arcor.de>
Subject: \captionsetup and keyval errmessages
Date: Sat, 20 Jan 2007 13:49:25 +0300

Hi Axel,

I occasionally mistyped option in \usepackage..{floatrow} and get error
message

----------------
! Package caption Error: wrongkey undefined.


See the caption package documentation for explanation.
Type  H <return>  for immediate help.
 ...

l.122 \ProcessOptionsWithKV{caption}
---------------


With other packages this error message addressed from keyval package itself.
I found that this error message ``corrected'' by \captionsetup command
inside caption3.sty. The same you could get in the case of wrong key in
hyperref after caption loaded.

My mind is in big mess: whether I can follow your ``hack'' to get correctly
addressed error messages or... I found that hyperref's keyval options get
error messages addressed from keyval unless the caption package loaded
before.
Maybe it is a good way to follow you and restore \KV@errx in
\AtEndOfPackage?

Best regards
Olga
\end{verbatim}

%******************************************************************************%

\bug{07-03-09}[v3.0k]{v3.0n}

\begin{verbatim}
To: caption@sommerfee.de
Message-Id: <B32FAB42-2B00-4DBA-91CD-C478E7C75262@gmx.de>
Subject: Bug in caption.sty
Date: Fri, 9 Mar 2007 18:33:47 +0100

Sehr geehrter Herr Sommerfeldt,

sie sind in der Beschreibung zum Caption-Style als Maintainer genannt.
Der aktuelle caption.sty (10.Jan.07, Teil der TexLive Distribution)  
scheint
Probleme zu haben, die Höhe umgebrochener Beschriftungen korrekt zu  
berechnen.
Der unten angefügte Testfall definiert zwei Minipages, die bündig zur  
unteren
Seite ausgerichtet sind. Die beiden angehängten PFDs zeigen die  
Ausgabe des
alten bzw. des neuen Styles. Meiner Meinung nach verhält sich der  
alte Style
korrekt. Gibt es für das Problem einen einfachen workaround?

\documentclass[a4paper,12pt,twoside,fleqn]{book}
\RequirePackage[font={small},labelfont= {small,bf,up},singlelinecheck=off,%
  tableposition=bottom]{caption}

\begin{document}
\framebox{
\begin{minipage}[b][2cm][b]{2cm}
bla bla
\end{minipage}}
\framebox{
\begin{minipage}[b][2cm][b]{2cm}
\captionof{figure}{Test}
\end{minipage}}

\end{document}
\end{verbatim}

\begin{example}{font={small},labelfont= {small,bf,up},singlelinecheck=off}
%\captionsetup[table]{position=bottom}
\vspace*{1ex}

\framebox{
\begin{minipage}[b][2cm][b]{2cm}
\small bla bla\strut
\end{minipage}}
\framebox{
\begin{minipage}[b][2cm][b]{2cm}
\captionof{figure}{Test}
\end{minipage}}

\vspace*{1ex}
\end{example}

%******************************************************************************%

\bug{07-09-13a}[v3.1]{v3.1a}

\begin{verbatim}
Date: Thu, 13 Sep 2007 09:43:52 +0200 (CEST)
Subject: Caption-Paket
To: caption@sommerfee.de
Message-ID: <23724.62814.qm@web26911.mail.ukl.yahoo.com>

Guten Tag Herr Sommerfeld,

bitte entschuldigen Sie, dass ich Sie direkt
anspreche.
Ist es möglich das die neueste Version Ihres
Caption-Paketes Probleme mit dem Listings-Paket
hervorruft?
Ich hatte gestern auf die neue Version umgestellt und
seit dem kann ich bei meinen Listings die Optionen
Captions und Label nicht mehr verwenden.
Kann ich irgendwie die alte Version wieder herstellen,
um weiterarbeiten zu können? Oder gibt es eine Art
Workaround?

Mit freundlichen Grüßen
Frank


Date: Thu, 13 Sep 2007 13:06:05 +0200 (CEST)
Subject: Re: Caption-Paket
To: Axel Sommerfeldt <caption@sommerfee.de>
Message-ID: <745787.93831.qm@web26906.mail.ukl.yahoo.com>

Hallo Herr Sommerfeld,

vielen Dank für Ihre schnelle Antwort.

> Könnten Sie mir ein Beispieldokument zuschicken,
> welches das Problem 
> aufzeigt? 
Bei diesen Einstellungen gab es das Problem. Ich habe
es auch schon auf das Caption-Paket eingegrenzt. Da
der Fehler verschwindet, sobald das Paket
auskommentiert wird.

\documentclass[a4paper, 12pt, german, halfparskip]{scrreprt} 

\usepackage[german]{babel}
\usepackage[T1]{fontenc}
\usepackage[latin1]{inputenc}

\usepackage{listings}

\usepackage[margin=10pt, font=small, labelfont=bf, format=plain]{caption, subfig}
\usepackage[pdfborder = false, bookmarksnumbered = true]{hyperref}

\begin{document}
  \begin{lstlisting}[float=htb, caption={Messung der Ausführungszeit},
                     label=list:ZeitMarkert]
    \\Code
  \end{lstlisting}
\end{document}

Das folgende ist ein Ausschnitt aus der log-Datei
! Argument of \strip@prefix has an extra }.
 
                \par 
l.22 ...Ausführungszeit}, label=list:ZeitMarkert ]
                                                  
I've run across a `}' that doesn't seem to match
anything.
For example, `\def\a#1{...}' and `\a}' would produce
this error. If you simply proceed now, the `\par' that
I've just inserted will cause me to report a runaway
argument that might be the root of the problem. But if
your `}' was spurious, just type `2' and it will go
away.

Runaway argument?
\relax
\end{verbatim}
%! Paragraph ended before \strip@prefix was complete.
% 
%                   \par 
%l.22 ...Ausführungszeit}, label=list:ZeitMarkert ]
%                                                  
%I suspect you've forgotten a `}', causing me to apply
%this
%control sequence to too much text. How can we recover?
%My plan is to forget the whole thing and hope for the
%best.
%
%Runaway argument?
%! Paragraph ended before \strip@period was complete.
% 
%                   \par 
%l.22 ...Ausführungszeit}, label=list:ZeitMarkert ]
%                                                  
%I suspect you've forgotten a `}', causing me to apply
%this
%control sequence to too much text. How can we recover?
%My plan is to forget the whole thing and hope for the
%best.
%
%! Undefined control sequence.
%\label ... \@currentlabelname \relax .\relax \@@@ 
%                                                 
%}\protected@write \@auxout...
%l.22 ...Ausführungszeit}, label=list:ZeitMarkert ]
%                                                  
%The control sequence at the end of the top line
%of your error message was never \def'ed. If you have
%misspelled it (e.g., `\hobx'), type `I' and the
%correct
%spelling (e.g., `I\hbox'). Otherwise just continue,
%and I'll forget about whatever was undefined.
\begin{verbatim}
Grüße
Frank
\end{verbatim}

\bigskip
\noindent\emph{Anmerkung: Der Fehlerfall wurde in das Testdokument
{\normalfont\texttt{hypcap\_base.tex}} eingepflegt}

%******************************************************************************%

\bug{07-09-13b}[v3.1]{v3.1a}

\begin{verbatim}
http://www.mrunix.de/forums/showthread.php?t=53862
Neues Caption Paket verträgt sich nicht mit Koma Klassen?

Hallo
Seit ich das neuste Caption Paket v3.1 mit den KOMA-Script package v.2.97a
benutze, bekomme ich die Fehlermeldung:

"Something's wrong -- perhaps a missing \item"

Dies nur dann, wenn ich beim Caption Paket die Optionen justification=centering,
labelsep=newline kombiniere und die Caption so lange ist, dass sie automatisch
auf mehrere Zeilen umbricht.

Hier mal ein Minimalbeispiel:
Code:

\documentclass[12pt, nochapterprefix, BCOR0mm]{scrreprt}
\usepackage[justification=centering,labelsep=newline]{caption}
\begin{document}
test
\begin{figure}
	\caption{sehr langer text sehr langer text sehr langer text sehr langer
      text sehr langer text sehr langer text}
\end{figure}			
\end{document}

Weiter scheint es so, dass beim Weglassen der Option justification=centering
zwar keine Fehlermeldung mehr erscheint, die Option labelsep=newline aber
unwirksam ist.

Ich möchte nun aber, dass das Label der Beschriftung auf einer Zeile steht und
darunter (im Idealfall zentriert) die Beschreibung.
Kann mir jemand weiterhelfen?

Besten Dank!
__________________
Grüsse

Thomas
\end{verbatim}

%******************************************************************************%

\bug{07-09-14}[v3.1]{v3.1a}

\begin{verbatim}
Date: Fri, 14 Sep 2007 01:45:53 +0200
To:  caption@sommerfee.de
Subject: [caption-Paket] Fehler mit neuem caption 3.1 und longtables

Hallo Axel,

ich habe gerade das neue caption-Paket installiert und kann nun meine
Dokumente nicht mehr texen. Das Problem sind captions in longtables:

\documentclass{scrartcl}
\usepackage{longtable}
\usepackage{caption}

\begin{document}
\begin{longtable}{c|c|c|c|c}
\caption[Longtable]{Longtable with caption}\\
\hline
1 & 2 & 3 & 4 & 5\tabularnewline
\hline
asd & s & s & s & asd\tabularnewline
\hline
\end{longtable}
\end{document}

Die Fehlermeldung lautet:
! You can't use a prefix with `\noalign'.
<to be read again>
                   \noalign
l.16 \caption
             [Longtable]{Longtable with caption}\\

Der Fehler taucht aber nur auf, wenn ich eine koma-skript Klasse verwende.
Gibt es da einen Workaround? caption for longtable zu laden hilft leider
nicht weiter. (Ich schreibe für das Programm LyX die Dokus und mit dem
neuen caption 3.1 kann man nun die Handbücher nicht mehr kompilieren.)

Danke und Gruß
Uwe
\end{verbatim}

\bigskip
\noindent\emph{Anmerkung: Der Fehlerfall wurde in das Testdokument
{\normalfont\texttt{koma-script.tex}} eingepflegt}

%******************************************************************************%

\bug{07-09-15}[v3.1]{v3.1a}

\begin{verbatim}
Date: Fri, 14 Sep 2007 15:48:18 +0200
To: Axel Sommerfeldt <caption@sommerfee.de>
Subject: Re: [caption-Paket] Fehler mit neuem caption 3.1 und longtables
In-Reply-To: <5.2.1.1.2.20070914064227.01a53c60@post.strato.de>

Axel Sommerfeldt schrieb:

Danke für den Bug-Report. Ich habe den Fehler nachgestellt, beseitigt,
und werde wohl in den nächsten Tagen eine neue Version 3.1a nach CTAN packen.

Wow, das ging aber schnell!

> 2.  Die dieser E-Mail anhängende Version 3.1a des caption-Paketes
> ausprobieren. Eine Rückmeldung, ob damit deine Probleme beseitigt sind,
> wäre toll!

Leider ist jetzt ein neuer Fehler drin:

-------
\documentclass{scrartcl}

\usepackage[T1]{fontenc}
\usepackage[latin9]{inputenc}

\usepackage{longtable}
\usepackage{subfigure}

\usepackage{caption}

\begin{document}

\begin{longtable}{c|c|c|c|c}
\caption[Longtable]{Longtable with caption}\\
\hline
1 & 2 & 3 & 4 & 5\tabularnewline
\hline
asd & s & s & s & asd\tabularnewline
\hline
\end{longtable}

\end{document}
-------

Die Fehlermeldung lautet nun:

! Undefined control sequence.
<argument> if\@captype
                       topcap
l.14 \caption[Longtable]{Longtable with caption}
                                                \\

Jetzt stört sich caption am Paket subfigure.

Danke und Gruß
Uwe
\end{verbatim}

%******************************************************************************%

\bug{07-09-17a}[v3.1]{v3.1b}

\begin{verbatim}
Subject: Caption / hyperref bug
Date: Mon, 17 Sep 2007 07:18:02 +0200
Message-ID: <9C48EEFF0A9A1147A920EBEF2CE8CBF0E1438D@STBEVS01.stb.sun.ac.za>
To: <caption@sommerfee.de>

Axel 

There is a clash between caption and hyperref in the newest versions of
both packages. It occurs with the breaklinks=true option of hyperref in
Latex. PdfLatex do not give the same error. See the example that follows

Danie Els 

   %--- test.tex ------------------- 
   \documentclass{article} 
   \usepackage{caption} 
   \usepackage{hyperref} 
   \hypersetup{breaklinks=true} %--> = false, then there is no error 
   \begin{document} 
      \begin{figure} 
      \end{figure} 
   \end{document} 
   %-------------------------------- 

If I run: 

   pdflatex test.tex  --> No errors 

   latex test.tex     --> This gives the following error: 

        ! You can't use `\prevdepth' in horizontal mode. 
            \caption@@anchor ...ption@hypcapspace }\prevdepth\@tempdima \endgroup 

This is for Miktex 2.7 Beta 2 

 *File List* 
 article.cls    2005/09/16 v1.4f Standard LaTeX document class 
  size10.clo    2005/09/16 v1.4f Standard LaTeX file (size option) 
 caption.sty    2007/09/16 v3.1a Customising captions (AR) 
caption3.sty    2007/09/16 v3.1a caption3 kernel (AR) 
  keyval.sty    1999/03/16 v1.13 key=value parser (DPC) 
hyperref.sty    2007/06/14 v6.76i Hypertext links for LaTeX 
 hycolor.sty    2007/04/11 v1.1 Code for color options of hyperref/bookmark (HO) 
  pd1enc.def    2007/06/14 v6.76i Hyperref: PDFDocEncoding definition (HO) 
etexcmds.sty    2007/05/06 v1.0 Providing prefix for e-TeX command names (HO) 
infwarerr.sty   2007/06/14 v1.1 Providing info/warning/message (HO) 
hyperref.cfg    2002/06/06 v1.2 hyperref configuration of TeXLive 
kvoptions.sty   2007/06/11 v2.7 Connects package keyval with LaTeX options (HO) 
     url.sty    2005/06/27  ver 3.2  Verb mode for urls, etc. 
  hdvips.def    2007/06/14 v6.76i Hyperref driver for dvips 
 pdfmark.def    2007/06/14 v6.76i Hyperref definitions for pdfmark specials 
 nameref.sty    2007/05/29 v2.31 Cross-referencing by name of section 
refcount.sty    2006/02/20 v3.0 Data extraction from references (HO)  
*********** 
\end{verbatim}

%******************************************************************************%

\bug{07-09-17b}[v3.1]{v3.1b}

\begin{verbatim}
http://www.mrunix.de/forums/showthread.php?t=53862
Neues Caption Paket verträgt sich nicht mit Koma Klassen?

Hi Axel
Hab zwar caption 3.1a noch nicht drauf, aber ich hab rausgefunden, woran es
vorher lag. Anscheinend verursacht \centering innerhalb der figure Umgebung
den Fehler. Hier ein Minimalbeispiel:

Code:

\documentclass{scrreprt}
\usepackage[format=plain,labelsep=newline]{caption}
\begin{document}
test
\begin{figure}
	\centering
	\caption{sehr langer text sehr langer text sehr langer text sehr langer
	    text sehr langer text sehr langer text}
\end{figure}			
\end{document}

Werde jetzt dann dein neues Packet drauf tun und schauen, ob sich was tut.

Grüsse

Thomas
\end{verbatim}

%******************************************************************************%

\bug{07-10-17}[v3.1]{v3.1c}

\begin{verbatim}
http://www.mrunix.de/forums/showthread.php?t=54472
Float too large Problem bei sidewaystable nach Miktex Update 

Hallo,

ich hab vor ein paar Tagen ein Miktex Update gemacht (weiss aber nicht mehr,
welche Pakete da alle aktualisiert wurden) und habe seit dem ein Problem mit
gedrehten Tabellen:

Ich nutze sidewaystable mit tabularx:

Code:

\begin{sidewaystable}[htb]
	\begin{tabularx}{\textwidth}{l|X|X}
        ...
	\end{tabularx}
	\caption{bla bla}
	\label{tab:konfiguebersicht}
\end{sidewaystable}

Jetzt spuckt mir LaTeX bei jeder gedrehten Tabelle die Warnung:
"Float too large for page by 1.0pt".

Am Inhalt der Tabelle scheint es nicht zu liegen. Auch wenn man in jede Spalte
nur einen Buchstaben einträgt kommt die Warnung. Und dass die Tabelle (oder
der Float) genau 1pt zu groß ist, ist ja auch etwas merkwürdig.

Ich hab nochmal ein kleines Beispiel gebastelt:

Code:

\documentclass[%
	pdftex,%        PDFTex verwenden
	a4paper,%       A4 Papier
]{scrreprt}

\usepackage{scrpage2}

\usepackage[german,ngerman]{babel}
\usepackage[latin1]{inputenc}
\usepackage[T1]{fontenc}

\usepackage[pdftex,
	]
	{hyperref}

\usepackage[
	bf,			% fett
	hang		% hängend	
] {caption}

\usepackage{tabularx}
\usepackage{rotating}


\begin{document}

\chapter{Hallo}

\begin{sidewaystable}[htb]
	\begin{tabularx}{\textwidth}{l|l|l}
		\textbf{A}	&	\textbf{B} 	&	\textbf{C}\\
	\end{tabularx}
	\caption{Caption}
	\label{tab:tabelle}
\end{sidewaystable}

\end{document}

Das komische ist: Die Warnung verschwindet, wenn man das caption- oder das
hyperref-Paket weglässt. Beide brauche ich aber.
\end{verbatim}

%******************************************************************************%

\bug{07-10-24}[v3.0f]{v3.1d}

\begin{verbatim}
http://www.mrunix.de/forums/showthread.php?t=54618
Package caption error: Undefined format 'normal'

Hi,

I'm writing my thesis in latex using Miktex as distribution and TeXnicCenter
as editor. I can compile my script but I'm getting every time the error
message mentioned above:
Code:

!Package caption error: Undefined format 'normal'
See the caption package documentation for explanation.

I just can't find out neither in the doc files nor over the net what this
kind of error suppose to mean. I can also get an output which seems quite
good to me.

I'm not sure if it's important ( but I learned, that it almost never not
important ) but I also get the warning:

Code:

Package caption warning: Unused \captionsetup[floatfoot] on input line 473.

This line is empty and lay between two subsections that are not written yet.

I downloaded Miktex yesterday and installed it this morning, so the reason
can't be an old package.

I looked into the caption Package directory and there are there three
*.sty files (caption.sty, caption2.sty, caption3.sty)

Has someone else this problem too?
Has anyone an idea where is this coming from?

THX

frymor
\end{verbatim}

%******************************************************************************%

\bug{07-11-09}[v3.0]{v3.1f}

\begin{verbatim}
Date: Fri, 09 Nov 2007 22:19:11 +0100
Newsgroups: comp.text.tex
Subject: figures, captions and skips
Message-ID: <4734cee5$0$13123$9b4e6d93@newsspool2.arcor-online.net>

Hi all,

while including figures/images of various widths, I noticed
that within the figure environment apparently the space/skip
generated between the caption and the included image changes
proportionally to the width of the image being included.  To
be precise, I would expect the code

\documentclass{book}

...

\usepackage{graphicx}
\usepackage{caption}
\captionsetup{position=bottom,font=small,labelfont=bf,format=hang}
\usepackage{subfig}

...

\begin{figure}
\centering
\includegraphics[keepaspectratio=true]{the-image.pdf}
\captionsetup{aboveskip=1em,belowskip=1em}
\caption[the title]{the caption}
\end{figure}

to place a caption below the image (which is assumed to be of
\textwidth width) using the same aboveskip amount that is being
generated for an image of half that width.  However, the code

...

\begin{figure}
\centering
\subfloat[][]{\includegraphics[keepaspectratio=true]{the-image.pdf}}
\captionsetup{aboveskip=1em,belowskip=1em}
\caption{...}
\end{figure}

...

does produce the intended aboveskip but with the drawback of a
subnumber on a one image figure which is somewhat undesirable.
Setting the global \abovecaptionskip had the same effect as
\captionsetup{aboveskip=1em}, ie, only producing an offset
bt not an absolute distance.

Is there a way to unconditionally/globally set the aboveskip
for a caption of a figure to a value that remains constant for
all images irrespective of their dimension ?

Thanks and regards,
Christian Boehme
\end{verbatim}
\bigskip
\begin{verbatim}
Date: Thu, 15 Nov 2007 03:18:51 +0100
Newsgroups: comp.text.tex
Subject: Re: figures, captions and skips
Message-ID: <473bac9d$0$27139$9b4e6d93@newsspool1.arcor-online.net>

Martin Heller wrote:

> \documentclass{book}
> \usepackage[demo]{graphicx}
> \usepackage{caption}
> \captionsetup{position=bottom,font=small,labelfont=bf,format=hang}
> \usepackage{subfig}
> \usepackage{lipsum}
> 
> \begin{document}
> 
> \lipsum[1]
> 
> \begin{figure}[hbt]
> \centering
> \includegraphics[keepaspectratio=true,width=\textwidth]{the-image.pdf}

Replace "width=\textwidth" with "width=0.5\textwidth" here and you
will notice that the spacing between the image and the caption decreases.

> \begin{figure}[hbt]
> \centering
> \subfloat[][]{\includegraphics[keepaspectratio=true,width=0.5\textwidth]

Doing the opposite here (ie, width=\textwidth) does _not_ change the
spacing, which was the point I was making.  However, it appears that the
caption package code introduces these variable spacings (why I have no
idea).  Throwing away all the \captionsetup{} statements in the figure
environments and a

\setlength{\abovecaptionskip}{0em}

in the preamble will work and show consistent behaviour but bombs me back
to where I started.  Needless to say, redefining \figurename does not
do anything.

Are there different Latex variants out there ?  I keep reading things
that just Don't Work (TM) when trying them out myself.  I use

pdfTeX using libpoppler 3.141592-1.40.3-2.2 (Web2C 7.5.6)

to generate PDFs from Latex.

Cheers,
Christian
\end{verbatim}
\clearpage

\begin{example}{}
\begin{figure}[H]
\centering
\includegraphics[keepaspectratio=true,width=0.5\linewidth]{the-image.pdf}
\captionsetup{aboveskip=1em,belowskip=1em}
\caption[the title]{the caption}
\end{figure}

\begin{figure}[H]
\centering
\includegraphics[keepaspectratio=true,width=0.7\linewidth]{the-image.pdf}
\captionsetup{aboveskip=1em,belowskip=1em}
\caption[the title]{the caption}
\end{figure}

\begin{figure}[H]
\centering
\includegraphics[keepaspectratio=true,width=1.0\linewidth]{the-image.pdf}
\captionsetup{aboveskip=1em,belowskip=1em}
\caption[the title]{the caption}
\end{figure}
\end{example}

%******************************************************************************%

\bug{07-12-06}[v3.1]{v3.1f}

\begin{verbatim}
Newsgroups: comp.text.tex
Subject: Re: Error using the caption package - Undefined control sequence
Date: Wed, 5 Dec 2007 05:37:26 -0800 (PST)
Message-ID: <b5bdef42-68c8-4b3e-8043-f635aaf92d28@v4g2000hsf.googlegroups.com>

On Dec 5, 1:05 pm, Ulrike Fischer <ne...@nililand.de> wrote:
> Am Wed, 5 Dec 2007 04:30:55 -0800 (PST) schrieb larse...@gmail.com:
>
> > Hi,
>
> > I just installed the newest version of the caption package as it was
> > required in order to use subfig. However, now I get an error message
> > upon using the \caption{} command which I, to my surprise, could not
> > find anywhere on the web. Applying pdflatex to the document
> > now results in the following message:
>
> > ! Undefined control sequence.
> > \caption@dblarg #1->\kernel@ifnextchar
>
> I can't reproduce the error but I doubt anyway that caption is at
> fault. \kernel@ifnextchar is here defined in latex.ltx (LaTeX2e
> <2005/12/01>):
>
> \documentclass{article}
> \makeatletter
> \show\kernel@ifnextchar
> \makeatother
>
> \begin{document}
> test
> \end{document}
>
> gives
>
> > \kernel@ifnextchar=\long macro:
>
> ....
> l.3 \show\kernel@ifnextchar
>
> Is it also defined on your side?
>
> --
> Ulrike Fischer

From the source, \kernel@ifnextchar was added 2004/01/23.  So any
LaTeX before that will be in trouble, I guess.

Joseph Wright
\end{verbatim}

%******************************************************************************%

\end{document}

